% =========================================================================
%  Figura «tipo Mathematica»: estado «disco cargado, resto neutros» del
%  sistema de cuatro conductores (combinación de los estados elementales con
%  phi = P e_C). Líneas de campo con flujo fijo por línea, equipotenciales y
%  signo de la carga inducida. Datos: figuras/datos/sistema_libre_*.
% =========================================================================
\begin{figure}[t]
\centering
\begin{tikzpicture}[x=0.66cm, y=0.66cm, font=\small]
  % Contornos de los cuatro conductores
% Generado por scripts/generar_datos.py (no editar a mano).
\def\formaA{(-4.250,-2.950) -- (-3.050,-3.950) -- (0.220,-0.180) -- cycle}
\def\formaB{(-2.750,1.250) -- (-0.650,1.250) -- (-0.650,2.850) -- (-2.750,2.850) -- cycle}
\def\formaC{(1.330,0.350) -- (1.330,0.363) -- (1.329,0.375) -- (1.329,0.388) -- (1.327,0.400) -- (1.326,0.413) -- (1.324,0.425) -- (1.322,0.437) -- (1.320,0.450) -- (1.317,0.462) -- (1.314,0.474) -- (1.310,0.486) -- (1.307,0.498) -- (1.302,0.510) -- (1.298,0.522) -- (1.293,0.534) -- (1.289,0.545) -- (1.283,0.557) -- (1.278,0.568) -- (1.272,0.579) -- (1.266,0.590) -- (1.259,0.601) -- (1.253,0.611) -- (1.246,0.622) -- (1.238,0.632) -- (1.231,0.642) -- (1.223,0.652) -- (1.215,0.662) -- (1.207,0.671) -- (1.198,0.680) -- (1.189,0.689) -- (1.180,0.698) -- (1.171,0.707) -- (1.162,0.715) -- (1.152,0.723) -- (1.142,0.731) -- (1.132,0.738) -- (1.122,0.746) -- (1.111,0.753) -- (1.101,0.759) -- (1.090,0.766) -- (1.079,0.772) -- (1.068,0.778) -- (1.057,0.783) -- (1.045,0.789) -- (1.034,0.793) -- (1.022,0.798) -- (1.010,0.802) -- (0.998,0.807) -- (0.986,0.810) -- (0.974,0.814) -- (0.962,0.817) -- (0.950,0.820) -- (0.937,0.822) -- (0.925,0.824) -- (0.913,0.826) -- (0.900,0.827) -- (0.888,0.829) -- (0.875,0.829) -- (0.863,0.830) -- (0.850,0.830) -- (0.837,0.830) -- (0.825,0.829) -- (0.812,0.829) -- (0.800,0.827) -- (0.787,0.826) -- (0.775,0.824) -- (0.763,0.822) -- (0.750,0.820) -- (0.738,0.817) -- (0.726,0.814) -- (0.714,0.810) -- (0.702,0.807) -- (0.690,0.802) -- (0.678,0.798) -- (0.666,0.793) -- (0.655,0.789) -- (0.643,0.783) -- (0.632,0.778) -- (0.621,0.772) -- (0.610,0.766) -- (0.599,0.759) -- (0.589,0.753) -- (0.578,0.746) -- (0.568,0.738) -- (0.558,0.731) -- (0.548,0.723) -- (0.538,0.715) -- (0.529,0.707) -- (0.520,0.698) -- (0.511,0.689) -- (0.502,0.680) -- (0.493,0.671) -- (0.485,0.662) -- (0.477,0.652) -- (0.469,0.642) -- (0.462,0.632) -- (0.454,0.622) -- (0.447,0.611) -- (0.441,0.601) -- (0.434,0.590) -- (0.428,0.579) -- (0.422,0.568) -- (0.417,0.557) -- (0.411,0.545) -- (0.407,0.534) -- (0.402,0.522) -- (0.398,0.510) -- (0.393,0.498) -- (0.390,0.486) -- (0.386,0.474) -- (0.383,0.462) -- (0.380,0.450) -- (0.378,0.437) -- (0.376,0.425) -- (0.374,0.413) -- (0.373,0.400) -- (0.371,0.388) -- (0.371,0.375) -- (0.370,0.363) -- (0.370,0.350) -- (0.370,0.337) -- (0.371,0.325) -- (0.371,0.312) -- (0.373,0.300) -- (0.374,0.287) -- (0.376,0.275) -- (0.378,0.263) -- (0.380,0.250) -- (0.383,0.238) -- (0.386,0.226) -- (0.390,0.214) -- (0.393,0.202) -- (0.398,0.190) -- (0.402,0.178) -- (0.407,0.166) -- (0.411,0.155) -- (0.417,0.143) -- (0.422,0.132) -- (0.428,0.121) -- (0.434,0.110) -- (0.441,0.099) -- (0.447,0.089) -- (0.454,0.078) -- (0.462,0.068) -- (0.469,0.058) -- (0.477,0.048) -- (0.485,0.038) -- (0.493,0.029) -- (0.502,0.020) -- (0.511,0.011) -- (0.520,0.002) -- (0.529,-0.007) -- (0.538,-0.015) -- (0.548,-0.023) -- (0.558,-0.031) -- (0.568,-0.038) -- (0.578,-0.046) -- (0.589,-0.053) -- (0.599,-0.059) -- (0.610,-0.066) -- (0.621,-0.072) -- (0.632,-0.078) -- (0.643,-0.083) -- (0.655,-0.089) -- (0.666,-0.093) -- (0.678,-0.098) -- (0.690,-0.102) -- (0.702,-0.107) -- (0.714,-0.110) -- (0.726,-0.114) -- (0.738,-0.117) -- (0.750,-0.120) -- (0.763,-0.122) -- (0.775,-0.124) -- (0.787,-0.126) -- (0.800,-0.127) -- (0.812,-0.129) -- (0.825,-0.129) -- (0.837,-0.130) -- (0.850,-0.130) -- (0.863,-0.130) -- (0.875,-0.129) -- (0.888,-0.129) -- (0.900,-0.127) -- (0.913,-0.126) -- (0.925,-0.124) -- (0.937,-0.122) -- (0.950,-0.120) -- (0.962,-0.117) -- (0.974,-0.114) -- (0.986,-0.110) -- (0.998,-0.107) -- (1.010,-0.102) -- (1.022,-0.098) -- (1.034,-0.093) -- (1.045,-0.089) -- (1.057,-0.083) -- (1.068,-0.078) -- (1.079,-0.072) -- (1.090,-0.066) -- (1.101,-0.059) -- (1.111,-0.053) -- (1.122,-0.046) -- (1.132,-0.038) -- (1.142,-0.031) -- (1.152,-0.023) -- (1.162,-0.015) -- (1.171,-0.007) -- (1.180,0.002) -- (1.189,0.011) -- (1.198,0.020) -- (1.207,0.029) -- (1.215,0.038) -- (1.223,0.048) -- (1.231,0.058) -- (1.238,0.068) -- (1.246,0.078) -- (1.253,0.089) -- (1.259,0.099) -- (1.266,0.110) -- (1.272,0.121) -- (1.278,0.132) -- (1.283,0.143) -- (1.289,0.155) -- (1.293,0.166) -- (1.298,0.178) -- (1.302,0.190) -- (1.307,0.202) -- (1.310,0.214) -- (1.314,0.226) -- (1.317,0.238) -- (1.320,0.250) -- (1.322,0.263) -- (1.324,0.275) -- (1.326,0.287) -- (1.327,0.300) -- (1.329,0.312) -- (1.329,0.325) -- (1.330,0.337) -- cycle}
\def\formaD{(1.050,1.320) -- (3.550,1.320) -- (3.550,1.680) -- (1.050,1.680) -- cycle}

  %% potenciales del estado libre (phi_C = 1)
\expandafter\def\csname philibreA\endcsname{0{,}23}
\expandafter\def\csname philibreB\endcsname{0{,}20}
\expandafter\def\csname philibreC\endcsname{1{,}00}
\expandafter\def\csname philibreD\endcsname{0{,}41}

  \fill[azulfondo!60!white] (-5,-4.5) rectangle (5,4.5);
  \begin{scope}
    \clip (-5,-4.5) rectangle (5,4.5);
    \tikzset{numEQ/.style={azulcielo, line width=0.35pt, densely dashed}}
    % Disco cargado, resto neutros - equipotenciales
% Generado por scripts/generar_datos.py (no editar a mano).
\EQ{(-3.260,-4.241) (-3.220,-4.246) (-3.180,-4.249) (-3.140,-4.252) (-3.100,-4.253) (-3.060,-4.254) (-3.020,-4.253) (-2.980,-4.251) (-2.940,-4.248) (-2.900,-4.244) (-2.860,-4.239) (-2.820,-4.233) (-2.780,-4.226) (-2.740,-4.219) (-2.700,-4.211) (-2.660,-4.203) (-2.620,-4.194) (-2.580,-4.185) (-2.540,-4.176) (-2.500,-4.166) (-2.460,-4.156) (-2.420,-4.146) (-2.380,-4.135) (-2.340,-4.125) (-2.300,-4.115) (-2.260,-4.104) (-2.220,-4.094) (-2.180,-4.083) (-2.140,-4.072) (-2.100,-4.062) (-2.060,-4.051) (-2.020,-4.041) (-1.980,-4.030) (-1.941,-4.020) (-1.900,-4.009) (-1.860,-3.999) (-1.820,-3.989) (-1.780,-3.979) (-1.740,-3.969) (-1.700,-3.959) (-1.660,-3.949) (-1.620,-3.940) (-1.580,-3.930) (-1.540,-3.921) (-1.500,-3.912) (-1.460,-3.903) (-1.420,-3.894) (-1.380,-3.886) (-1.340,-3.877) (-1.300,-3.869) (-1.260,-3.861) (-1.220,-3.853) (-1.180,-3.846) (-1.140,-3.838) (-1.100,-3.831) (-1.060,-3.824) (-1.020,-3.818) (-0.980,-3.811) (-0.940,-3.805) (-0.900,-3.799) (-0.860,-3.793) (-0.820,-3.787) (-0.780,-3.782) (-0.740,-3.777) (-0.700,-3.772) (-0.660,-3.767) (-0.620,-3.763) (-0.580,-3.759) (-0.540,-3.755) (-0.500,-3.751) (-0.460,-3.747) (-0.420,-3.744) (-0.380,-3.741) (-0.340,-3.737) (-0.300,-3.735) (-0.260,-3.732) (-0.220,-3.729) (-0.180,-3.727) (-0.140,-3.724) (-0.100,-3.722) (-0.060,-3.720) (-0.020,-3.718) (0.020,-3.716) (0.060,-3.715) (0.100,-3.713) (0.140,-3.711) (0.180,-3.710) (0.220,-3.708) (0.260,-3.707) (0.300,-3.706) (0.340,-3.704) (0.380,-3.703) (0.420,-3.702) (0.460,-3.700) (0.500,-3.699) (0.540,-3.698) (0.580,-3.696) (0.620,-3.695) (0.660,-3.693) (0.700,-3.692) (0.740,-3.690) (0.780,-3.689) (0.820,-3.687) (0.860,-3.685) (0.900,-3.684) (0.940,-3.682) (0.980,-3.680) (1.020,-3.678) (1.060,-3.675) (1.100,-3.673) (1.140,-3.670) (1.180,-3.668) (1.220,-3.665) (1.260,-3.662) (1.300,-3.659) (1.340,-3.656) (1.380,-3.653) (1.420,-3.649) (1.460,-3.645) (1.500,-3.641) (1.540,-3.637) (1.580,-3.633) (1.620,-3.628) (1.660,-3.624) (1.700,-3.619) (1.740,-3.613) (1.780,-3.608) (1.820,-3.602) (1.860,-3.596) (1.900,-3.590) (1.940,-3.584) (1.980,-3.577) (2.020,-3.570) (2.060,-3.562) (2.100,-3.555) (2.140,-3.547) (2.180,-3.538) (2.220,-3.529) (2.260,-3.520) (2.300,-3.511) (2.340,-3.501) (2.380,-3.491) (2.420,-3.480) (2.460,-3.469) (2.500,-3.457) (2.540,-3.445) (2.580,-3.433) (2.618,-3.420) (2.660,-3.406) (2.700,-3.391) (2.740,-3.377) (2.780,-3.361) (2.820,-3.345) (2.860,-3.328) (2.900,-3.310) (2.940,-3.292) (2.980,-3.273) (3.020,-3.253) (3.060,-3.232) (3.100,-3.210) (3.140,-3.187) (3.180,-3.163) (3.217,-3.140) (3.260,-3.112) (3.300,-3.084) (3.334,-3.060) (3.380,-3.025) (3.412,-3.000) (3.460,-2.960) (3.500,-2.924) (3.540,-2.887) (3.580,-2.848) (3.620,-2.806) (3.660,-2.762) (3.696,-2.720) (3.729,-2.680) (3.760,-2.640) (3.790,-2.600) (3.818,-2.560) (3.846,-2.520) (3.872,-2.480) (3.896,-2.440) (3.920,-2.401) (3.940,-2.366) (3.960,-2.330) (3.980,-2.293) (4.000,-2.254) (4.020,-2.214) (4.040,-2.172) (4.060,-2.129) (4.080,-2.084) (4.099,-2.040) (4.115,-2.000) (4.131,-1.960) (4.146,-1.920) (4.161,-1.880) (4.175,-1.840) (4.189,-1.800) (4.203,-1.760) (4.216,-1.720) (4.228,-1.680) (4.241,-1.640) (4.253,-1.600) (4.264,-1.560) (4.275,-1.520) (4.286,-1.480) (4.297,-1.440) (4.307,-1.400) (4.317,-1.360) (4.327,-1.320) (4.336,-1.280) (4.345,-1.240) (4.354,-1.200) (4.363,-1.160) (4.371,-1.120) (4.379,-1.080) (4.387,-1.040) (4.395,-1.000) (4.403,-0.960) (4.410,-0.920) (4.417,-0.880) (4.424,-0.840) (4.431,-0.800) (4.438,-0.760) (4.444,-0.720) (4.451,-0.680) (4.457,-0.640) (4.463,-0.600) (4.469,-0.560) (4.475,-0.520) (4.481,-0.480) (4.486,-0.440) (4.492,-0.400) (4.497,-0.360) (4.502,-0.320) (4.507,-0.280) (4.512,-0.240) (4.517,-0.200) (4.522,-0.160) (4.527,-0.120) (4.532,-0.080) (4.536,-0.040) (4.541,-0.000) (4.545,0.040) (4.550,0.080) (4.554,0.120) (4.558,0.160) (4.562,0.200) (4.566,0.240) (4.570,0.280) (4.574,0.320) (4.578,0.360) (4.582,0.400) (4.585,0.440) (4.589,0.480) (4.592,0.520) (4.596,0.560) (4.599,0.600) (4.602,0.640) (4.605,0.680) (4.608,0.720) (4.611,0.760) (4.613,0.800) (4.616,0.840) (4.618,0.880) (4.620,0.920) (4.622,0.960) (4.624,1.000) (4.626,1.040) (4.627,1.080) (4.628,1.120) (4.629,1.160) (4.630,1.200) (4.630,1.240) (4.631,1.280) (4.630,1.320) (4.630,1.360) (4.629,1.400) (4.628,1.440) (4.627,1.480) (4.625,1.520) (4.623,1.560) (4.621,1.600) (4.618,1.640) (4.615,1.680) (4.611,1.720) (4.608,1.760) (4.603,1.800) (4.599,1.840) (4.593,1.880) (4.588,1.920) (4.582,1.960) (4.575,2.000) (4.569,2.040) (4.561,2.080) (4.553,2.120) (4.545,2.160) (4.536,2.200) (4.527,2.240) (4.517,2.280) (4.507,2.320) (4.496,2.360) (4.484,2.400) (4.472,2.440) (4.459,2.480) (4.445,2.520) (4.431,2.560) (4.416,2.600) (4.400,2.640) (4.384,2.680) (4.366,2.720) (4.348,2.760) (4.329,2.800) (4.308,2.840) (4.287,2.880) (4.264,2.920) (4.240,2.960) (4.214,3.000) (4.187,3.040) (4.160,3.078) (4.128,3.120) (4.100,3.154) (4.060,3.200) (4.023,3.240) (3.982,3.280) (3.940,3.319) (3.900,3.353) (3.866,3.380) (3.820,3.415) (3.783,3.440) (3.740,3.468) (3.700,3.493) (3.660,3.516) (3.620,3.537) (3.580,3.558) (3.540,3.577) (3.500,3.595) (3.460,3.612) (3.420,3.628) (3.380,3.643) (3.340,3.657) (3.300,3.671) (3.260,3.684) (3.220,3.696) (3.180,3.707) (3.140,3.718) (3.100,3.728) (3.060,3.738) (3.020,3.748) (2.980,3.756) (2.940,3.765) (2.900,3.773) (2.861,3.780) (2.820,3.787) (2.780,3.794) (2.740,3.800) (2.700,3.806) (2.660,3.812) (2.620,3.818) (2.580,3.823) (2.540,3.828) (2.500,3.832) (2.460,3.837) (2.420,3.841) (2.380,3.845) (2.340,3.849) (2.300,3.852) (2.260,3.855) (2.220,3.858) (2.180,3.861) (2.140,3.864) (2.100,3.867) (2.060,3.869) (2.020,3.871) (1.980,3.873) (1.940,3.875) (1.900,3.877) (1.860,3.878) (1.820,3.880) (1.780,3.881) (1.740,3.882) (1.700,3.883) (1.660,3.884) (1.620,3.885) (1.580,3.886) (1.540,3.886) (1.500,3.887) (1.460,3.887) (1.420,3.888) (1.380,3.888) (1.340,3.888) (1.300,3.888) (1.260,3.888) (1.220,3.887) (1.180,3.887) (1.140,3.886) (1.100,3.886) (1.060,3.885) (1.020,3.885) (0.980,3.884) (0.940,3.883) (0.900,3.882) (0.860,3.881) (0.820,3.880) (0.780,3.878) (0.740,3.877) (0.700,3.876) (0.660,3.874) (0.620,3.873) (0.580,3.871) (0.540,3.870) (0.500,3.868) (0.460,3.866) (0.420,3.865) (0.380,3.863) (0.340,3.861) (0.300,3.859) (0.260,3.858) (0.220,3.856) (0.180,3.854) (0.140,3.852) (0.100,3.850) (0.060,3.848) (0.020,3.846) (-0.020,3.845) (-0.060,3.843) (-0.100,3.841) (-0.140,3.840) (-0.180,3.838) (-0.220,3.836) (-0.260,3.835) (-0.300,3.833) (-0.340,3.832) (-0.380,3.830) (-0.420,3.829) (-0.460,3.828) (-0.500,3.827) (-0.540,3.826) (-0.580,3.825) (-0.620,3.824) (-0.660,3.823) (-0.700,3.822) (-0.740,3.821) (-0.780,3.820) (-0.820,3.819) (-0.860,3.818) (-0.900,3.817) (-0.940,3.817) (-0.980,3.816) (-1.020,3.815) (-1.060,3.814) (-1.100,3.813) (-1.140,3.813) (-1.180,3.812) (-1.220,3.811) (-1.260,3.810) (-1.300,3.809) (-1.340,3.808) (-1.380,3.807) (-1.420,3.806) (-1.460,3.804) (-1.500,3.803) (-1.540,3.802) (-1.580,3.800) (-1.620,3.798) (-1.660,3.797) (-1.700,3.795) (-1.740,3.793) (-1.780,3.791) (-1.820,3.788) (-1.860,3.786) (-1.900,3.783) (-1.940,3.781) (-1.980,3.777) (-2.020,3.774) (-2.060,3.771) (-2.100,3.767) (-2.140,3.763) (-2.180,3.758) (-2.220,3.754) (-2.260,3.749) (-2.300,3.743) (-2.340,3.737) (-2.380,3.731) (-2.420,3.725) (-2.460,3.717) (-2.500,3.710) (-2.540,3.701) (-2.580,3.693) (-2.620,3.683) (-2.660,3.673) (-2.700,3.662) (-2.740,3.650) (-2.780,3.638) (-2.820,3.624) (-2.860,3.610) (-2.900,3.595) (-2.940,3.578) (-2.980,3.561) (-3.020,3.542) (-3.060,3.522) (-3.100,3.501) (-3.136,3.480) (-3.180,3.453) (-3.220,3.426) (-3.257,3.400) (-3.300,3.368) (-3.334,3.340) (-3.379,3.300) (-3.420,3.261) (-3.460,3.220) (-3.496,3.180) (-3.529,3.140) (-3.560,3.100) (-3.589,3.060) (-3.616,3.020) (-3.640,2.982) (-3.665,2.940) (-3.687,2.900) (-3.708,2.860) (-3.727,2.820) (-3.745,2.780) (-3.762,2.740) (-3.779,2.700) (-3.794,2.660) (-3.808,2.620) (-3.821,2.580) (-3.834,2.540) (-3.846,2.500) (-3.857,2.460) (-3.868,2.420) (-3.877,2.380) (-3.887,2.340) (-3.895,2.300) (-3.903,2.260) (-3.911,2.220) (-3.918,2.180) (-3.925,2.140) (-3.931,2.100) (-3.937,2.060) (-3.942,2.020) (-3.948,1.980) (-3.952,1.940) (-3.957,1.900) (-3.961,1.860) (-3.964,1.820) (-3.968,1.780) (-3.971,1.740) (-3.974,1.700) (-3.977,1.660) (-3.979,1.620) (-3.981,1.580) (-3.983,1.540) (-3.985,1.500) (-3.987,1.460) (-3.988,1.420) (-3.990,1.380) (-3.991,1.340) (-3.992,1.300) (-3.993,1.260) (-3.994,1.220) (-3.994,1.180) (-3.995,1.140) (-3.996,1.100) (-3.996,1.060) (-3.997,1.020) (-3.997,0.980) (-3.998,0.940) (-3.999,0.900) (-3.999,0.860) (-4.000,0.820) (-4.000,0.780) (-4.001,0.740) (-4.002,0.700) (-4.003,0.660) (-4.004,0.620) (-4.005,0.580) (-4.006,0.540) (-4.007,0.500) (-4.009,0.460) (-4.010,0.420) (-4.012,0.380) (-4.014,0.340) (-4.016,0.300) (-4.018,0.260) (-4.020,0.220) (-4.023,0.180) (-4.026,0.140) (-4.029,0.100) (-4.032,0.060) (-4.035,0.020) (-4.039,-0.020) (-4.042,-0.060) (-4.046,-0.100) (-4.050,-0.140) (-4.055,-0.180) (-4.059,-0.220) (-4.064,-0.260) (-4.069,-0.300) (-4.074,-0.340) (-4.080,-0.380) (-4.085,-0.420) (-4.091,-0.460) (-4.097,-0.500) (-4.104,-0.540) (-4.110,-0.580) (-4.117,-0.620) (-4.124,-0.660) (-4.131,-0.700) (-4.139,-0.740) (-4.146,-0.780) (-4.154,-0.820) (-4.162,-0.860) (-4.170,-0.900) (-4.179,-0.940) (-4.187,-0.980) (-4.196,-1.020) (-4.205,-1.060) (-4.215,-1.100) (-4.224,-1.140) (-4.234,-1.180) (-4.243,-1.220) (-4.253,-1.260) (-4.264,-1.300) (-4.274,-1.340) (-4.284,-1.380) (-4.295,-1.420) (-4.306,-1.460) (-4.316,-1.500) (-4.327,-1.540) (-4.338,-1.580) (-4.350,-1.620) (-4.361,-1.660) (-4.372,-1.700) (-4.383,-1.740) (-4.395,-1.780) (-4.406,-1.820) (-4.418,-1.860) (-4.429,-1.900) (-4.441,-1.940) (-4.453,-1.980) (-4.464,-2.020) (-4.475,-2.060) (-4.487,-2.100) (-4.498,-2.140) (-4.509,-2.180) (-4.520,-2.220) (-4.531,-2.260) (-4.541,-2.300) (-4.552,-2.340) (-4.562,-2.380) (-4.571,-2.420) (-4.581,-2.460) (-4.590,-2.500) (-4.598,-2.540) (-4.606,-2.580) (-4.613,-2.620) (-4.620,-2.660) (-4.626,-2.700) (-4.631,-2.740) (-4.635,-2.780) (-4.638,-2.820) (-4.641,-2.860) (-4.642,-2.900) (-4.642,-2.940) (-4.641,-2.980) (-4.639,-3.020) (-4.636,-3.060) (-4.632,-3.100) (-4.626,-3.140) (-4.620,-3.180) (-4.612,-3.220) (-4.603,-3.260) (-4.592,-3.300) (-4.581,-3.340) (-4.568,-3.380) (-4.553,-3.420) (-4.538,-3.460) (-4.520,-3.500) (-4.501,-3.540) (-4.480,-3.580) (-4.460,-3.616) (-4.432,-3.660) (-4.404,-3.700) (-4.373,-3.740) (-4.340,-3.779) (-4.300,-3.820) (-4.260,-3.857) (-4.220,-3.891) (-4.181,-3.920) (-4.140,-3.949) (-4.100,-3.974) (-4.060,-3.997) (-4.020,-4.018) (-3.980,-4.038) (-3.940,-4.056) (-3.900,-4.073) (-3.860,-4.089) (-3.820,-4.104) (-3.780,-4.119) (-3.740,-4.132) (-3.700,-4.145) (-3.660,-4.157) (-3.620,-4.168) (-3.580,-4.178) (-3.540,-4.189) (-3.500,-4.198) (-3.460,-4.207) (-3.420,-4.215) (-3.380,-4.223) (-3.340,-4.230) (-3.300,-4.236) (-3.260,-4.241)}
\EQ{(-3.120,-4.042) (-3.080,-4.047) (-3.040,-4.046) (-3.000,-4.041) (-2.960,-4.031) (-2.920,-4.019) (-2.880,-4.003) (-2.840,-3.985) (-2.800,-3.966) (-2.760,-3.945) (-2.720,-3.923) (-2.680,-3.900) (-2.645,-3.880) (-2.600,-3.853) (-2.560,-3.829) (-2.520,-3.805) (-2.480,-3.781) (-2.440,-3.756) (-2.400,-3.731) (-2.360,-3.707) (-2.320,-3.682) (-2.285,-3.660) (-2.240,-3.632) (-2.200,-3.608) (-2.160,-3.583) (-2.122,-3.560) (-2.080,-3.534) (-2.040,-3.510) (-2.000,-3.486) (-1.960,-3.462) (-1.922,-3.440) (-1.880,-3.415) (-1.840,-3.392) (-1.800,-3.370) (-1.760,-3.347) (-1.720,-3.325) (-1.680,-3.303) (-1.640,-3.282) (-1.600,-3.261) (-1.560,-3.240) (-1.520,-3.220) (-1.480,-3.201) (-1.440,-3.182) (-1.400,-3.163) (-1.360,-3.145) (-1.320,-3.127) (-1.280,-3.110) (-1.240,-3.094) (-1.200,-3.078) (-1.160,-3.063) (-1.120,-3.049) (-1.080,-3.035) (-1.040,-3.022) (-1.000,-3.010) (-0.960,-2.998) (-0.920,-2.987) (-0.880,-2.977) (-0.840,-2.967) (-0.800,-2.958) (-0.760,-2.950) (-0.720,-2.942) (-0.680,-2.935) (-0.640,-2.929) (-0.600,-2.923) (-0.560,-2.918) (-0.520,-2.913) (-0.480,-2.909) (-0.440,-2.906) (-0.400,-2.903) (-0.360,-2.901) (-0.320,-2.899) (-0.280,-2.897) (-0.240,-2.896) (-0.200,-2.895) (-0.160,-2.894) (-0.120,-2.894) (-0.080,-2.894) (-0.040,-2.894) (-0.000,-2.895) (0.040,-2.895) (0.080,-2.896) (0.120,-2.897) (0.160,-2.898) (0.200,-2.899) (0.240,-2.900) (0.280,-2.901) (0.320,-2.902) (0.360,-2.903) (0.400,-2.904) (0.440,-2.905) (0.480,-2.905) (0.520,-2.906) (0.560,-2.907) (0.600,-2.907) (0.640,-2.907) (0.680,-2.907) (0.720,-2.907) (0.760,-2.907) (0.800,-2.907) (0.840,-2.906) (0.880,-2.905) (0.920,-2.904) (0.960,-2.902) (1.000,-2.901) (1.040,-2.899) (1.080,-2.896) (1.120,-2.894) (1.160,-2.891) (1.200,-2.888) (1.240,-2.884) (1.280,-2.880) (1.320,-2.876) (1.360,-2.871) (1.400,-2.866) (1.440,-2.861) (1.480,-2.855) (1.520,-2.849) (1.560,-2.843) (1.600,-2.836) (1.640,-2.828) (1.680,-2.820) (1.720,-2.812) (1.760,-2.803) (1.800,-2.794) (1.840,-2.784) (1.880,-2.774) (1.920,-2.763) (1.960,-2.752) (2.000,-2.740) (2.040,-2.728) (2.080,-2.715) (2.120,-2.701) (2.160,-2.687) (2.200,-2.672) (2.240,-2.657) (2.280,-2.640) (2.320,-2.624) (2.360,-2.606) (2.400,-2.588) (2.440,-2.568) (2.480,-2.548) (2.520,-2.527) (2.560,-2.506) (2.600,-2.483) (2.639,-2.460) (2.680,-2.435) (2.720,-2.409) (2.760,-2.382) (2.800,-2.355) (2.840,-2.326) (2.874,-2.300) (2.920,-2.264) (2.960,-2.231) (2.995,-2.200) (3.040,-2.160) (3.080,-2.122) (3.120,-2.083) (3.160,-2.042) (3.199,-2.000) (3.234,-1.960) (3.268,-1.920) (3.300,-1.881) (3.332,-1.840) (3.360,-1.803) (3.391,-1.760) (3.420,-1.720) (3.447,-1.680) (3.473,-1.640) (3.498,-1.600) (3.520,-1.563) (3.546,-1.520) (3.568,-1.480) (3.590,-1.440) (3.611,-1.400) (3.632,-1.360) (3.652,-1.320) (3.671,-1.280) (3.690,-1.240) (3.708,-1.200) (3.726,-1.160) (3.743,-1.120) (3.760,-1.080) (3.776,-1.040) (3.791,-1.000) (3.807,-0.960) (3.821,-0.920) (3.836,-0.880) (3.850,-0.840) (3.863,-0.800) (3.877,-0.760) (3.890,-0.720) (3.902,-0.680) (3.915,-0.640) (3.927,-0.600) (3.938,-0.560) (3.950,-0.520) (3.961,-0.480) (3.972,-0.440) (3.983,-0.400) (3.993,-0.360) (4.004,-0.320) (4.014,-0.280) (4.024,-0.240) (4.033,-0.200) (4.043,-0.160) (4.053,-0.120) (4.062,-0.080) (4.071,-0.040) (4.080,-0.000) (4.089,0.040) (4.098,0.080) (4.106,0.120) (4.115,0.160) (4.123,0.200) (4.132,0.240) (4.140,0.280) (4.148,0.320) (4.156,0.360) (4.164,0.400) (4.172,0.440) (4.180,0.480) (4.187,0.520) (4.195,0.560) (4.202,0.600) (4.209,0.640) (4.216,0.680) (4.223,0.720) (4.230,0.760) (4.236,0.800) (4.242,0.840) (4.248,0.880) (4.254,0.920) (4.259,0.960) (4.263,1.000) (4.268,1.040) (4.272,1.080) (4.275,1.120) (4.278,1.160) (4.281,1.200) (4.282,1.240) (4.284,1.280) (4.284,1.320) (4.284,1.360) (4.284,1.400) (4.282,1.440) (4.280,1.480) (4.277,1.520) (4.274,1.560) (4.269,1.600) (4.264,1.640) (4.258,1.680) (4.251,1.720) (4.243,1.760) (4.234,1.800) (4.224,1.840) (4.213,1.880) (4.201,1.920) (4.188,1.960) (4.174,2.000) (4.160,2.038) (4.143,2.080) (4.125,2.120) (4.106,2.160) (4.086,2.200) (4.065,2.240) (4.042,2.280) (4.020,2.316) (3.992,2.360) (3.964,2.400) (3.940,2.433) (3.904,2.480) (3.871,2.520) (3.840,2.555) (3.800,2.598) (3.760,2.638) (3.720,2.675) (3.680,2.710) (3.644,2.740) (3.600,2.775) (3.566,2.800) (3.520,2.832) (3.480,2.859) (3.446,2.880) (3.400,2.907) (3.360,2.930) (3.320,2.951) (3.280,2.971) (3.240,2.991) (3.200,3.009) (3.160,3.026) (3.120,3.042) (3.080,3.058) (3.040,3.073) (3.000,3.087) (2.961,3.100) (2.920,3.113) (2.880,3.125) (2.840,3.137) (2.800,3.148) (2.760,3.158) (2.720,3.168) (2.680,3.177) (2.640,3.186) (2.600,3.195) (2.560,3.203) (2.520,3.210) (2.480,3.217) (2.440,3.224) (2.400,3.231) (2.360,3.237) (2.320,3.242) (2.280,3.248) (2.240,3.253) (2.200,3.257) (2.160,3.262) (2.120,3.266) (2.080,3.270) (2.040,3.274) (2.000,3.277) (1.960,3.280) (1.920,3.283) (1.880,3.285) (1.840,3.288) (1.800,3.290) (1.760,3.292) (1.720,3.293) (1.680,3.295) (1.640,3.296) (1.600,3.297) (1.560,3.298) (1.520,3.298) (1.480,3.299) (1.440,3.299) (1.400,3.299) (1.360,3.299) (1.320,3.298) (1.280,3.297) (1.240,3.297) (1.200,3.296) (1.160,3.294) (1.120,3.293) (1.080,3.291) (1.040,3.290) (1.000,3.287) (0.960,3.285) (0.920,3.283) (0.880,3.280) (0.840,3.277) (0.800,3.274) (0.760,3.271) (0.720,3.268) (0.680,3.264) (0.640,3.260) (0.600,3.256) (0.560,3.252) (0.520,3.248) (0.480,3.243) (0.440,3.238) (0.400,3.234) (0.360,3.228) (0.320,3.223) (0.280,3.218) (0.240,3.212) (0.200,3.207) (0.160,3.201) (0.120,3.195) (0.080,3.189) (0.040,3.184) (-0.000,3.178) (-0.040,3.172) (-0.080,3.166) (-0.120,3.161) (-0.160,3.156) (-0.200,3.151) (-0.240,3.146) (-0.280,3.142) (-0.320,3.138) (-0.360,3.135) (-0.400,3.132) (-0.440,3.130) (-0.480,3.129) (-0.520,3.128) (-0.560,3.128) (-0.600,3.129) (-0.640,3.129) (-0.680,3.130) (-0.720,3.131) (-0.760,3.132) (-0.800,3.133) (-0.840,3.133) (-0.880,3.133) (-0.920,3.134) (-0.960,3.134) (-1.000,3.134) (-1.040,3.134) (-1.080,3.134) (-1.120,3.134) (-1.160,3.133) (-1.200,3.133) (-1.240,3.133) (-1.280,3.132) (-1.320,3.132) (-1.360,3.131) (-1.400,3.131) (-1.440,3.130) (-1.480,3.129) (-1.520,3.129) (-1.560,3.128) (-1.600,3.127) (-1.640,3.126) (-1.680,3.125) (-1.720,3.124) (-1.760,3.122) (-1.800,3.121) (-1.840,3.120) (-1.880,3.118) (-1.920,3.116) (-1.960,3.114) (-2.000,3.112) (-2.040,3.110) (-2.080,3.107) (-2.120,3.105) (-2.160,3.102) (-2.200,3.099) (-2.240,3.095) (-2.280,3.091) (-2.320,3.087) (-2.360,3.082) (-2.400,3.077) (-2.440,3.071) (-2.480,3.064) (-2.520,3.056) (-2.560,3.048) (-2.600,3.038) (-2.640,3.026) (-2.680,3.013) (-2.720,2.996) (-2.760,2.975) (-2.800,2.949) (-2.833,2.920) (-2.860,2.890) (-2.895,2.840) (-2.916,2.800) (-2.934,2.760) (-2.950,2.720) (-2.964,2.680) (-2.976,2.640) (-2.987,2.600) (-2.997,2.560) (-3.006,2.520) (-3.014,2.480) (-3.022,2.440) (-3.028,2.400) (-3.035,2.360) (-3.041,2.320) (-3.046,2.280) (-3.051,2.240) (-3.055,2.200) (-3.059,2.160) (-3.063,2.120) (-3.067,2.080) (-3.070,2.040) (-3.072,2.000) (-3.075,1.960) (-3.077,1.920) (-3.078,1.880) (-3.080,1.840) (-3.081,1.800) (-3.082,1.760) (-3.082,1.720) (-3.082,1.680) (-3.082,1.640) (-3.081,1.600) (-3.080,1.560) (-3.078,1.520) (-3.076,1.480) (-3.073,1.440) (-3.069,1.400) (-3.065,1.360) (-3.060,1.320) (-3.054,1.280) (-3.047,1.240) (-3.039,1.200) (-3.030,1.160) (-3.021,1.120) (-3.011,1.080) (-3.000,1.040) (-2.990,1.000) (-2.980,0.960) (-2.969,0.920) (-2.960,0.880) (-2.950,0.840) (-2.942,0.800) (-2.934,0.760) (-2.927,0.720) (-2.921,0.680) (-2.915,0.640) (-2.911,0.600) (-2.907,0.560) (-2.903,0.520) (-2.901,0.480) (-2.899,0.440) (-2.898,0.400) (-2.897,0.360) (-2.897,0.320) (-2.898,0.280) (-2.899,0.240) (-2.901,0.200) (-2.903,0.160) (-2.906,0.120) (-2.909,0.080) (-2.913,0.040) (-2.917,-0.000) (-2.921,-0.040) (-2.927,-0.080) (-2.932,-0.120) (-2.938,-0.160) (-2.945,-0.200) (-2.952,-0.240) (-2.960,-0.280) (-2.968,-0.320) (-2.976,-0.360) (-2.985,-0.400) (-2.995,-0.440) (-3.005,-0.480) (-3.016,-0.520) (-3.028,-0.560) (-3.040,-0.600) (-3.052,-0.640) (-3.065,-0.680) (-3.079,-0.720) (-3.093,-0.760) (-3.108,-0.800) (-3.124,-0.840) (-3.140,-0.880) (-3.156,-0.920) (-3.174,-0.960) (-3.192,-1.000) (-3.210,-1.040) (-3.229,-1.080) (-3.249,-1.120) (-3.269,-1.160) (-3.290,-1.200) (-3.312,-1.240) (-3.334,-1.280) (-3.356,-1.320) (-3.379,-1.360) (-3.400,-1.396) (-3.427,-1.440) (-3.451,-1.480) (-3.476,-1.520) (-3.500,-1.558) (-3.527,-1.600) (-3.553,-1.640) (-3.580,-1.680) (-3.607,-1.720) (-3.634,-1.760) (-3.660,-1.798) (-3.689,-1.840) (-3.718,-1.880) (-3.746,-1.920) (-3.775,-1.960) (-3.800,-1.995) (-3.832,-2.040) (-3.860,-2.078) (-3.890,-2.120) (-3.919,-2.160) (-3.948,-2.200) (-3.978,-2.240) (-4.007,-2.280) (-4.035,-2.320) (-4.060,-2.354) (-4.092,-2.400) (-4.120,-2.439) (-4.148,-2.480) (-4.175,-2.520) (-4.200,-2.558) (-4.226,-2.600) (-4.250,-2.640) (-4.273,-2.680) (-4.294,-2.720) (-4.313,-2.760) (-4.329,-2.800) (-4.342,-2.840) (-4.352,-2.880) (-4.358,-2.920) (-4.359,-2.960) (-4.356,-3.000) (-4.348,-3.040) (-4.337,-3.080) (-4.322,-3.120) (-4.304,-3.160) (-4.283,-3.200) (-4.260,-3.238) (-4.233,-3.280) (-4.204,-3.320) (-4.173,-3.360) (-4.140,-3.399) (-4.104,-3.440) (-4.066,-3.480) (-4.025,-3.520) (-3.981,-3.560) (-3.940,-3.596) (-3.900,-3.629) (-3.860,-3.660) (-3.820,-3.690) (-3.780,-3.718) (-3.740,-3.746) (-3.700,-3.772) (-3.660,-3.797) (-3.623,-3.820) (-3.580,-3.845) (-3.540,-3.868) (-3.500,-3.890) (-3.460,-3.911) (-3.420,-3.931) (-3.380,-3.950) (-3.340,-3.969) (-3.300,-3.986) (-3.260,-4.002) (-3.220,-4.017) (-3.180,-4.029) (-3.140,-4.039) (-3.120,-4.042)}
\EQ{(-0.047,-0.507) (-0.072,-0.540) (-0.100,-0.573) (-0.138,-0.620) (-0.160,-0.654) (-0.195,-0.700) (-0.220,-0.732) (-0.257,-0.780) (-0.280,-0.814) (-0.312,-0.860) (-0.340,-0.900) (-0.367,-0.940) (-0.392,-0.980) (-0.416,-1.020) (-0.440,-1.059) (-0.460,-1.094) (-0.480,-1.133) (-0.500,-1.175) (-0.519,-1.220) (-0.533,-1.260) (-0.545,-1.300) (-0.553,-1.340) (-0.557,-1.380) (-0.557,-1.420) (-0.553,-1.460) (-0.544,-1.500) (-0.530,-1.540) (-0.511,-1.580) (-0.486,-1.620) (-0.460,-1.655) (-0.420,-1.700) (-0.380,-1.740) (-0.340,-1.774) (-0.308,-1.800) (-0.260,-1.835) (-0.224,-1.860) (-0.180,-1.888) (-0.140,-1.912) (-0.100,-1.934) (-0.060,-1.955) (-0.020,-1.975) (0.020,-1.994) (0.060,-2.012) (0.100,-2.029) (0.140,-2.045) (0.178,-2.060) (0.220,-2.075) (0.260,-2.089) (0.300,-2.102) (0.340,-2.114) (0.380,-2.126) (0.420,-2.136) (0.460,-2.146) (0.500,-2.156) (0.540,-2.165) (0.580,-2.173) (0.620,-2.180) (0.660,-2.187) (0.700,-2.193) (0.740,-2.199) (0.780,-2.203) (0.820,-2.208) (0.860,-2.211) (0.900,-2.215) (0.940,-2.217) (0.980,-2.219) (1.020,-2.220) (1.060,-2.221) (1.100,-2.221) (1.140,-2.220) (1.180,-2.219) (1.220,-2.217) (1.260,-2.215) (1.300,-2.212) (1.340,-2.208) (1.380,-2.204) (1.420,-2.199) (1.460,-2.193) (1.500,-2.187) (1.540,-2.180) (1.580,-2.172) (1.620,-2.164) (1.660,-2.155) (1.700,-2.145) (1.740,-2.135) (1.780,-2.124) (1.820,-2.112) (1.860,-2.099) (1.900,-2.086) (1.940,-2.072) (1.980,-2.057) (2.020,-2.041) (2.060,-2.024) (2.100,-2.006) (2.140,-1.988) (2.180,-1.969) (2.220,-1.948) (2.260,-1.927) (2.300,-1.904) (2.340,-1.881) (2.374,-1.860) (2.420,-1.831) (2.460,-1.804) (2.494,-1.780) (2.540,-1.746) (2.574,-1.720) (2.620,-1.683) (2.660,-1.649) (2.693,-1.620) (2.737,-1.580) (2.778,-1.540) (2.818,-1.500) (2.856,-1.460) (2.892,-1.420) (2.920,-1.387) (2.959,-1.340) (2.991,-1.300) (3.020,-1.262) (3.051,-1.220) (3.080,-1.180) (3.107,-1.140) (3.133,-1.100) (3.158,-1.060) (3.180,-1.025) (3.207,-0.980) (3.229,-0.940) (3.252,-0.900) (3.273,-0.860) (3.294,-0.820) (3.314,-0.780) (3.333,-0.740) (3.352,-0.700) (3.370,-0.660) (3.388,-0.620) (3.405,-0.580) (3.422,-0.540) (3.438,-0.500) (3.454,-0.460) (3.470,-0.420) (3.485,-0.380) (3.499,-0.340) (3.514,-0.300) (3.528,-0.260) (3.542,-0.220) (3.555,-0.180) (3.569,-0.140) (3.582,-0.100) (3.595,-0.060) (3.608,-0.020) (3.620,0.019) (3.633,0.060) (3.645,0.100) (3.658,0.140) (3.670,0.180) (3.682,0.220) (3.695,0.260) (3.707,0.300) (3.719,0.340) (3.732,0.380) (3.744,0.420) (3.757,0.460) (3.770,0.500) (3.782,0.540) (3.795,0.580) (3.808,0.620) (3.820,0.658) (3.833,0.700) (3.846,0.740) (3.859,0.780) (3.871,0.820) (3.884,0.860) (3.896,0.900) (3.907,0.940) (3.918,0.980) (3.929,1.020) (3.939,1.060) (3.948,1.100) (3.956,1.140) (3.964,1.180) (3.970,1.220) (3.975,1.260) (3.979,1.300) (3.982,1.340) (3.984,1.380) (3.984,1.420) (3.983,1.460) (3.981,1.500) (3.977,1.540) (3.972,1.580) (3.965,1.620) (3.956,1.660) (3.946,1.700) (3.934,1.740) (3.920,1.779) (3.903,1.820) (3.885,1.860) (3.864,1.900) (3.841,1.940) (3.815,1.980) (3.787,2.020) (3.760,2.054) (3.720,2.100) (3.682,2.140) (3.640,2.180) (3.600,2.215) (3.569,2.240) (3.520,2.277) (3.487,2.300) (3.440,2.331) (3.400,2.356) (3.360,2.379) (3.322,2.400) (3.283,2.420) (3.243,2.440) (3.200,2.460) (3.160,2.477) (3.120,2.494) (3.080,2.510) (3.040,2.525) (3.000,2.539) (2.960,2.553) (2.920,2.566) (2.880,2.578) (2.840,2.590) (2.800,2.601) (2.760,2.611) (2.720,2.621) (2.680,2.631) (2.640,2.640) (2.600,2.648) (2.560,2.656) (2.520,2.664) (2.480,2.671) (2.440,2.678) (2.400,2.684) (2.360,2.690) (2.320,2.696) (2.280,2.702) (2.240,2.707) (2.200,2.711) (2.160,2.716) (2.120,2.720) (2.080,2.724) (2.040,2.727) (2.000,2.731) (1.960,2.734) (1.920,2.737) (1.880,2.739) (1.840,2.741) (1.800,2.743) (1.760,2.745) (1.720,2.747) (1.680,2.748) (1.640,2.749) (1.600,2.750) (1.560,2.750) (1.520,2.750) (1.480,2.750) (1.440,2.750) (1.400,2.750) (1.360,2.749) (1.320,2.748) (1.280,2.747) (1.240,2.745) (1.200,2.744) (1.160,2.741) (1.120,2.739) (1.080,2.736) (1.040,2.733) (1.000,2.730) (0.960,2.727) (0.920,2.723) (0.880,2.718) (0.840,2.713) (0.800,2.708) (0.760,2.703) (0.720,2.697) (0.680,2.690) (0.640,2.683) (0.600,2.676) (0.560,2.668) (0.520,2.659) (0.480,2.650) (0.440,2.640) (0.400,2.630) (0.360,2.619) (0.320,2.607) (0.280,2.594) (0.240,2.580) (0.200,2.565) (0.160,2.549) (0.120,2.532) (0.080,2.514) (0.040,2.494) (-0.000,2.472) (-0.040,2.448) (-0.080,2.422) (-0.120,2.392) (-0.159,2.360) (-0.200,2.321) (-0.238,2.280) (-0.270,2.240) (-0.299,2.200) (-0.320,2.166) (-0.340,2.131) (-0.360,2.093) (-0.380,2.050) (-0.400,2.003) (-0.416,1.960) (-0.430,1.920) (-0.443,1.880) (-0.455,1.840) (-0.467,1.800) (-0.477,1.760) (-0.487,1.720) (-0.497,1.680) (-0.506,1.640) (-0.514,1.600) (-0.523,1.560) (-0.531,1.520) (-0.539,1.480) (-0.548,1.440) (-0.556,1.400) (-0.566,1.360) (-0.577,1.320) (-0.590,1.280) (-0.610,1.240) (-0.640,1.206) (-0.677,1.180) (-0.713,1.160) (-0.753,1.140) (-0.794,1.120) (-0.836,1.100) (-0.877,1.080) (-0.916,1.060) (-0.954,1.040) (-0.989,1.020) (-1.040,0.989) (-1.080,0.963) (-1.113,0.940) (-1.160,0.904) (-1.200,0.870) (-1.232,0.840) (-1.260,0.811) (-1.288,0.780) (-1.320,0.740) (-1.349,0.700) (-1.374,0.660) (-1.397,0.620) (-1.416,0.580) (-1.434,0.540) (-1.449,0.500) (-1.462,0.460) (-1.473,0.420) (-1.483,0.380) (-1.491,0.340) (-1.497,0.300) (-1.501,0.260) (-1.505,0.220) (-1.507,0.180) (-1.507,0.140) (-1.507,0.100) (-1.505,0.060) (-1.502,0.020) (-1.497,-0.020) (-1.491,-0.060) (-1.484,-0.100) (-1.476,-0.140) (-1.466,-0.180) (-1.455,-0.220) (-1.443,-0.260) (-1.429,-0.300) (-1.413,-0.340) (-1.395,-0.380) (-1.375,-0.420) (-1.352,-0.460) (-1.325,-0.500) (-1.300,-0.534) (-1.260,-0.578) (-1.220,-0.614) (-1.183,-0.640) (-1.146,-0.660) (-1.100,-0.679) (-1.060,-0.690) (-1.020,-0.696) (-0.980,-0.699) (-0.940,-0.698) (-0.900,-0.695) (-0.860,-0.689) (-0.820,-0.681) (-0.780,-0.671) (-0.740,-0.660) (-0.700,-0.648) (-0.660,-0.634) (-0.620,-0.619) (-0.580,-0.602) (-0.540,-0.585) (-0.500,-0.568) (-0.460,-0.550) (-0.420,-0.530) (-0.380,-0.510) (-0.340,-0.492) (-0.300,-0.472) (-0.260,-0.451) (-0.220,-0.432) (-0.180,-0.409) (-0.140,-0.388) (-0.108,-0.368)}
\EQ{(0.021,-0.419) (-0.015,-0.460) (-0.040,-0.496)}
\EQ{(0.235,-0.185) (0.225,-0.240) (0.210,-0.280) (0.196,-0.320) (0.182,-0.360) (0.169,-0.400) (0.156,-0.440) (0.145,-0.480) (0.134,-0.520) (0.124,-0.560) (0.116,-0.600) (0.109,-0.640) (0.103,-0.680) (0.099,-0.720) (0.096,-0.760) (0.094,-0.800) (0.094,-0.840) (0.096,-0.880) (0.100,-0.920) (0.106,-0.960) (0.114,-1.000) (0.124,-1.040) (0.136,-1.080) (0.151,-1.120) (0.168,-1.160) (0.189,-1.200) (0.213,-1.240) (0.240,-1.279) (0.272,-1.320) (0.300,-1.351) (0.328,-1.380) (0.360,-1.409) (0.397,-1.440) (0.440,-1.472) (0.480,-1.499) (0.516,-1.520) (0.552,-1.540) (0.592,-1.560) (0.637,-1.580) (0.680,-1.597) (0.720,-1.612) (0.760,-1.625) (0.800,-1.636) (0.840,-1.646) (0.880,-1.655) (0.920,-1.663) (0.960,-1.670) (1.000,-1.675) (1.040,-1.680) (1.080,-1.683) (1.120,-1.686) (1.160,-1.687) (1.200,-1.687) (1.240,-1.687) (1.280,-1.685) (1.320,-1.683) (1.360,-1.679) (1.400,-1.675) (1.440,-1.669) (1.480,-1.663) (1.520,-1.656) (1.560,-1.648) (1.600,-1.638) (1.640,-1.628) (1.680,-1.617) (1.720,-1.605) (1.760,-1.592) (1.800,-1.578) (1.840,-1.562) (1.880,-1.546) (1.920,-1.529) (1.960,-1.510) (2.000,-1.490) (2.040,-1.469) (2.080,-1.447) (2.120,-1.424) (2.158,-1.400) (2.200,-1.373) (2.240,-1.345) (2.275,-1.320) (2.320,-1.285) (2.351,-1.260) (2.399,-1.220) (2.440,-1.183) (2.480,-1.145) (2.520,-1.105) (2.560,-1.062) (2.598,-1.020) (2.632,-0.980) (2.660,-0.946) (2.696,-0.900) (2.720,-0.867) (2.753,-0.820) (2.780,-0.781) (2.807,-0.740) (2.832,-0.700) (2.856,-0.660) (2.879,-0.620) (2.900,-0.581) (2.920,-0.544) (2.940,-0.505) (2.960,-0.465) (2.980,-0.423) (3.000,-0.380) (3.017,-0.340) (3.035,-0.300) (3.051,-0.260) (3.067,-0.220) (3.083,-0.180) (3.098,-0.140) (3.113,-0.100) (3.127,-0.060) (3.141,-0.020) (3.155,0.020) (3.169,0.060) (3.182,0.100) (3.195,0.140) (3.208,0.180) (3.222,0.220) (3.235,0.260) (3.248,0.300) (3.261,0.340) (3.275,0.380) (3.289,0.420) (3.303,0.460) (3.318,0.500) (3.333,0.540) (3.350,0.580) (3.366,0.620) (3.384,0.660) (3.403,0.700) (3.422,0.740) (3.443,0.780) (3.465,0.820) (3.488,0.860) (3.512,0.900) (3.537,0.940) (3.560,0.976) (3.588,1.020) (3.614,1.060) (3.638,1.100) (3.660,1.138) (3.680,1.176) (3.700,1.220) (3.715,1.260) (3.727,1.300) (3.736,1.340) (3.742,1.380) (3.746,1.420) (3.747,1.460) (3.746,1.500) (3.742,1.540) (3.737,1.580) (3.728,1.620) (3.717,1.660) (3.701,1.700) (3.681,1.740) (3.655,1.780) (3.621,1.820) (3.580,1.858) (3.540,1.889) (3.500,1.915) (3.460,1.938) (3.420,1.959) (3.380,1.977) (3.340,1.995) (3.300,2.011) (3.260,2.025) (3.220,2.039) (3.180,2.052) (3.140,2.064) (3.100,2.075) (3.060,2.086) (3.020,2.096) (2.980,2.106) (2.940,2.115) (2.900,2.124) (2.860,2.132) (2.820,2.139) (2.780,2.147) (2.740,2.154) (2.700,2.160) (2.660,2.166) (2.620,2.172) (2.580,2.178) (2.540,2.183) (2.500,2.188) (2.460,2.193) (2.420,2.197) (2.380,2.201) (2.340,2.205) (2.300,2.209) (2.260,2.213) (2.220,2.216) (2.180,2.219) (2.140,2.222) (2.100,2.225) (2.060,2.227) (2.020,2.229) (1.980,2.231) (1.940,2.233) (1.900,2.235) (1.860,2.237) (1.820,2.238) (1.780,2.239) (1.740,2.240) (1.700,2.241) (1.660,2.242) (1.620,2.243) (1.580,2.243) (1.540,2.243) (1.500,2.243) (1.460,2.243) (1.420,2.243) (1.380,2.243) (1.340,2.242) (1.300,2.241) (1.260,2.240) (1.220,2.239) (1.180,2.238) (1.140,2.236) (1.100,2.234) (1.060,2.232) (1.020,2.230) (0.980,2.227) (0.940,2.224) (0.900,2.220) (0.860,2.216) (0.820,2.212) (0.780,2.207) (0.740,2.201) (0.700,2.195) (0.660,2.188) (0.621,2.180) (0.580,2.171) (0.540,2.161) (0.500,2.151) (0.460,2.139) (0.420,2.126) (0.380,2.111) (0.340,2.095) (0.300,2.078) (0.263,2.060) (0.224,2.040) (0.189,2.020) (0.140,1.990) (0.100,1.962) (0.060,1.932) (0.022,1.900) (-0.020,1.862) (-0.060,1.822) (-0.097,1.780) (-0.130,1.740) (-0.160,1.701) (-0.189,1.660) (-0.215,1.620) (-0.240,1.580) (-0.264,1.540) (-0.286,1.500) (-0.307,1.460) (-0.328,1.420) (-0.348,1.380) (-0.368,1.340) (-0.389,1.300) (-0.410,1.260) (-0.431,1.220) (-0.454,1.180) (-0.479,1.140) (-0.505,1.100) (-0.533,1.060) (-0.560,1.022) (-0.591,0.980) (-0.619,0.940) (-0.647,0.900) (-0.673,0.860) (-0.696,0.820) (-0.718,0.780) (-0.738,0.740) (-0.755,0.700) (-0.770,0.660) (-0.783,0.620) (-0.793,0.580) (-0.801,0.540) (-0.807,0.500) (-0.811,0.460) (-0.812,0.420) (-0.811,0.380) (-0.808,0.340) (-0.801,0.300) (-0.793,0.260) (-0.781,0.220) (-0.767,0.180) (-0.749,0.140) (-0.727,0.100) (-0.701,0.060) (-0.670,0.020) (-0.640,-0.012) (-0.610,-0.040) (-0.560,-0.079) (-0.520,-0.104) (-0.480,-0.125) (-0.440,-0.143) (-0.400,-0.157) (-0.360,-0.169) (-0.320,-0.178) (-0.280,-0.186) (-0.240,-0.191) (-0.200,-0.194) (-0.160,-0.196) (-0.120,-0.196) (-0.080,-0.194) (-0.040,-0.192) (-0.000,-0.188) (0.040,-0.183) (0.080,-0.178) (0.120,-0.172) (0.160,-0.166) (0.200,-0.162) (0.216,-0.164)}
\EQ{(3.560,1.342) (3.565,1.400) (3.567,1.440) (3.567,1.480) (3.567,1.520) (3.566,1.560) (3.564,1.600) (3.560,1.640) (3.560,1.641)}
\EQ{(3.513,1.700) (3.460,1.709) (3.420,1.713) (3.380,1.717) (3.340,1.720) (3.300,1.723) (3.260,1.725) (3.220,1.728) (3.180,1.730) (3.140,1.732) (3.100,1.734) (3.060,1.735) (3.020,1.737) (2.980,1.738) (2.940,1.740) (2.900,1.741) (2.860,1.743) (2.820,1.744) (2.780,1.745) (2.740,1.746) (2.700,1.747) (2.660,1.748) (2.620,1.749) (2.580,1.750) (2.540,1.751) (2.500,1.751) (2.460,1.752) (2.420,1.753) (2.380,1.754) (2.340,1.754) (2.300,1.755) (2.260,1.755) (2.220,1.756) (2.180,1.756) (2.140,1.757) (2.100,1.757) (2.060,1.758) (2.020,1.758) (1.980,1.758) (1.940,1.759) (1.900,1.759) (1.860,1.759) (1.820,1.759) (1.780,1.760) (1.740,1.760) (1.700,1.760) (1.660,1.760) (1.620,1.760) (1.580,1.760) (1.540,1.760) (1.500,1.760) (1.460,1.760) (1.420,1.760) (1.380,1.760) (1.340,1.760) (1.300,1.761) (1.260,1.761) (1.220,1.761) (1.180,1.761) (1.140,1.762) (1.100,1.762) (1.060,1.764) (1.020,1.767) (0.980,1.772) (0.940,1.777) (0.900,1.783) (0.860,1.788) (0.820,1.792) (0.780,1.795) (0.740,1.797) (0.700,1.797) (0.660,1.795) (0.620,1.792) (0.580,1.787) (0.540,1.780) (0.500,1.772) (0.460,1.762) (0.420,1.750) (0.380,1.736) (0.340,1.720) (0.300,1.702) (0.260,1.682) (0.222,1.660) (0.180,1.634) (0.140,1.606) (0.106,1.580) (0.060,1.542) (0.020,1.505) (-0.020,1.464) (-0.059,1.420) (-0.093,1.380) (-0.120,1.345) (-0.153,1.300) (-0.180,1.261) (-0.207,1.220) (-0.232,1.180) (-0.257,1.140) (-0.280,1.101) (-0.304,1.060) (-0.326,1.020) (-0.348,0.980) (-0.368,0.940) (-0.387,0.900) (-0.405,0.860) (-0.422,0.820) (-0.437,0.780) (-0.450,0.740) (-0.461,0.700) (-0.471,0.660) (-0.478,0.620) (-0.483,0.580) (-0.486,0.540) (-0.487,0.500) (-0.485,0.460) (-0.481,0.420) (-0.474,0.380) (-0.464,0.340) (-0.452,0.300) (-0.436,0.260) (-0.416,0.220) (-0.392,0.180) (-0.364,0.140) (-0.329,0.100) (-0.300,0.071) (-0.263,0.040) (-0.220,0.009) (-0.180,-0.015) (-0.140,-0.035) (-0.100,-0.052) (-0.060,-0.067) (-0.020,-0.079) (0.020,-0.090) (0.060,-0.099) (0.100,-0.107) (0.140,-0.114) (0.180,-0.123) (0.220,-0.135) (0.260,-0.178) (0.269,-0.220) (0.270,-0.260) (0.269,-0.300) (0.268,-0.340) (0.268,-0.380) (0.268,-0.420) (0.270,-0.460) (0.272,-0.500) (0.277,-0.540) (0.282,-0.580) (0.290,-0.620) (0.299,-0.660) (0.310,-0.700) (0.324,-0.740) (0.340,-0.780) (0.358,-0.820) (0.380,-0.860) (0.405,-0.900) (0.433,-0.940) (0.460,-0.974) (0.500,-1.018) (0.540,-1.055) (0.580,-1.088) (0.620,-1.117) (0.655,-1.140) (0.700,-1.166) (0.740,-1.186) (0.780,-1.204) (0.820,-1.220) (0.860,-1.233) (0.900,-1.245) (0.940,-1.255) (0.980,-1.264) (1.020,-1.271) (1.060,-1.277) (1.100,-1.281) (1.140,-1.283) (1.180,-1.285) (1.220,-1.285) (1.260,-1.284) (1.300,-1.281) (1.340,-1.278) (1.380,-1.273) (1.420,-1.266) (1.460,-1.259) (1.500,-1.250) (1.540,-1.240) (1.580,-1.229) (1.620,-1.216) (1.660,-1.203) (1.700,-1.187) (1.740,-1.171) (1.780,-1.153) (1.820,-1.133) (1.860,-1.113) (1.900,-1.090) (1.940,-1.066) (1.980,-1.040) (2.020,-1.013) (2.060,-0.984) (2.100,-0.952) (2.139,-0.920) (2.180,-0.883) (2.220,-0.845) (2.260,-0.805) (2.300,-0.761) (2.335,-0.720) (2.368,-0.680) (2.398,-0.640) (2.427,-0.600) (2.455,-0.560) (2.480,-0.521) (2.505,-0.480) (2.528,-0.440) (2.550,-0.400) (2.571,-0.360) (2.591,-0.320) (2.610,-0.280) (2.628,-0.240) (2.645,-0.200) (2.660,-0.161) (2.676,-0.120) (2.690,-0.080) (2.703,-0.040) (2.716,-0.000) (2.728,0.040) (2.740,0.080) (2.751,0.120) (2.761,0.160) (2.770,0.200) (2.779,0.240) (2.788,0.280) (2.796,0.320) (2.804,0.360) (2.811,0.400) (2.818,0.440) (2.825,0.480) (2.832,0.520) (2.838,0.560) (2.844,0.600) (2.850,0.640) (2.857,0.680) (2.863,0.720) (2.870,0.760) (2.878,0.800) (2.886,0.840) (2.895,0.880) (2.906,0.920) (2.919,0.960) (2.935,1.000) (2.956,1.040) (2.980,1.077) (3.019,1.120) (3.060,1.152) (3.100,1.176) (3.140,1.195) (3.180,1.210) (3.220,1.224) (3.260,1.235) (3.300,1.245) (3.340,1.254) (3.380,1.263) (3.420,1.272) (3.460,1.280) (3.500,1.290) (3.530,1.300)}
\EQ{(0.980,-0.903) (1.020,-0.909) (1.060,-0.914) (1.100,-0.917) (1.140,-0.918) (1.180,-0.918) (1.220,-0.916) (1.260,-0.912) (1.300,-0.907) (1.340,-0.900) (1.380,-0.891) (1.420,-0.881) (1.460,-0.869) (1.500,-0.856) (1.540,-0.840) (1.580,-0.823) (1.620,-0.804) (1.660,-0.784) (1.700,-0.761) (1.740,-0.736) (1.780,-0.708) (1.818,-0.680) (1.860,-0.646) (1.900,-0.611) (1.932,-0.580) (1.971,-0.540) (2.000,-0.508) (2.040,-0.460) (2.070,-0.420) (2.098,-0.380) (2.120,-0.346) (2.148,-0.300) (2.170,-0.260) (2.190,-0.220) (2.209,-0.180) (2.226,-0.140) (2.242,-0.100) (2.256,-0.060) (2.269,-0.020) (2.281,0.020) (2.291,0.060) (2.300,0.100) (2.308,0.140) (2.314,0.180) (2.319,0.220) (2.323,0.260) (2.325,0.300) (2.326,0.340) (2.325,0.380) (2.323,0.420) (2.319,0.460) (2.314,0.500) (2.307,0.540) (2.298,0.580) (2.287,0.620) (2.273,0.660) (2.257,0.700) (2.238,0.740) (2.215,0.780) (2.189,0.820) (2.160,0.857) (2.121,0.900) (2.080,0.937) (2.040,0.969) (2.000,0.997) (1.962,1.020) (1.926,1.040) (1.886,1.060) (1.842,1.080) (1.800,1.097) (1.760,1.111) (1.720,1.125) (1.680,1.137) (1.640,1.148) (1.600,1.159) (1.560,1.168) (1.520,1.177) (1.480,1.185) (1.440,1.193) (1.400,1.200) (1.360,1.206) (1.320,1.213) (1.280,1.219) (1.240,1.224) (1.200,1.230) (1.160,1.237) (1.120,1.244) (1.080,1.253) (1.040,1.269) (1.000,1.295) (0.960,1.330) (0.925,1.360) (0.880,1.394) (0.840,1.417) (0.800,1.435) (0.760,1.448) (0.720,1.457) (0.680,1.463) (0.640,1.465) (0.600,1.465) (0.560,1.461) (0.520,1.455) (0.480,1.447) (0.440,1.436) (0.400,1.422) (0.360,1.406) (0.320,1.387) (0.280,1.365) (0.240,1.340) (0.200,1.312) (0.160,1.281) (0.120,1.246) (0.080,1.207) (0.040,1.164) (0.004,1.120) (-0.027,1.080) (-0.054,1.040) (-0.080,1.000) (-0.103,0.960) (-0.125,0.920) (-0.144,0.880) (-0.160,0.843) (-0.177,0.800) (-0.190,0.760) (-0.201,0.720) (-0.211,0.680) (-0.218,0.640) (-0.223,0.600) (-0.226,0.560) (-0.226,0.520) (-0.225,0.480) (-0.220,0.440) (-0.213,0.400) (-0.203,0.360) (-0.190,0.320) (-0.174,0.280) (-0.154,0.240) (-0.130,0.200) (-0.101,0.160) (-0.067,0.120) (-0.025,0.080) (0.020,0.043) (0.060,0.015) (0.100,-0.010) (0.140,-0.033) (0.180,-0.056) (0.220,-0.080) (0.260,-0.111) (0.299,-0.160) (0.319,-0.200) (0.333,-0.240) (0.346,-0.280) (0.358,-0.320) (0.371,-0.360) (0.386,-0.400) (0.402,-0.440) (0.420,-0.480) (0.440,-0.518) (0.465,-0.560) (0.492,-0.600) (0.520,-0.637) (0.559,-0.680) (0.600,-0.720) (0.640,-0.753) (0.678,-0.780) (0.720,-0.807) (0.760,-0.828) (0.800,-0.847) (0.840,-0.864) (0.880,-0.877) (0.920,-0.889) (0.960,-0.899) (0.980,-0.903)}
\EQ{(0.920,-0.640) (0.960,-0.648) (1.000,-0.653) (1.040,-0.656) (1.080,-0.657) (1.120,-0.655) (1.160,-0.652) (1.200,-0.647) (1.239,-0.640) (1.280,-0.631) (1.320,-0.619) (1.360,-0.606) (1.400,-0.591) (1.440,-0.573) (1.480,-0.553) (1.520,-0.530) (1.560,-0.505) (1.595,-0.480) (1.640,-0.444) (1.680,-0.409) (1.710,-0.380) (1.740,-0.348) (1.780,-0.300) (1.810,-0.260) (1.837,-0.220) (1.860,-0.183) (1.880,-0.147) (1.900,-0.107) (1.920,-0.061) (1.936,-0.020) (1.949,0.020) (1.961,0.060) (1.970,0.100) (1.978,0.140) (1.984,0.180) (1.988,0.220) (1.991,0.260) (1.991,0.300) (1.990,0.340) (1.987,0.380) (1.981,0.420) (1.974,0.460) (1.964,0.500) (1.952,0.540) (1.938,0.580) (1.920,0.620) (1.900,0.660) (1.877,0.700) (1.849,0.740) (1.820,0.777) (1.782,0.820) (1.740,0.860) (1.700,0.893) (1.664,0.920) (1.620,0.950) (1.580,0.974) (1.540,0.995) (1.500,1.015) (1.460,1.033) (1.420,1.050) (1.380,1.065) (1.340,1.079) (1.300,1.092) (1.260,1.104) (1.220,1.115) (1.180,1.126) (1.140,1.137) (1.100,1.148) (1.061,1.160) (1.020,1.173) (0.980,1.187) (0.940,1.202) (0.900,1.216) (0.860,1.229) (0.820,1.240) (0.780,1.248) (0.740,1.254) (0.700,1.257) (0.660,1.258) (0.620,1.255) (0.580,1.250) (0.540,1.242) (0.500,1.231) (0.460,1.218) (0.420,1.201) (0.380,1.181) (0.343,1.160) (0.300,1.132) (0.260,1.102) (0.220,1.067) (0.180,1.027) (0.140,0.981) (0.108,0.940) (0.081,0.900) (0.060,0.864) (0.040,0.827) (0.020,0.784) (0.003,0.740) (-0.011,0.700) (-0.021,0.660) (-0.030,0.620) (-0.036,0.580) (-0.039,0.540) (-0.040,0.500) (-0.038,0.460) (-0.034,0.420) (-0.027,0.380) (-0.017,0.340) (-0.004,0.300) (0.012,0.260) (0.032,0.220) (0.056,0.180) (0.080,0.146) (0.118,0.100) (0.156,0.060) (0.198,0.020) (0.240,-0.018) (0.280,-0.056) (0.320,-0.100) (0.350,-0.140) (0.376,-0.180) (0.400,-0.220) (0.424,-0.260) (0.448,-0.300) (0.475,-0.340) (0.500,-0.374) (0.539,-0.420) (0.577,-0.460) (0.620,-0.498) (0.660,-0.528) (0.700,-0.554) (0.740,-0.576) (0.780,-0.595) (0.820,-0.612) (0.860,-0.625) (0.900,-0.636) (0.920,-0.640)}
\EQ{(0.940,-0.462) (0.980,-0.465) (1.020,-0.465) (1.060,-0.463) (1.100,-0.459) (1.140,-0.452) (1.180,-0.443) (1.220,-0.431) (1.260,-0.417) (1.300,-0.400) (1.340,-0.380) (1.375,-0.360) (1.420,-0.331) (1.460,-0.301) (1.500,-0.266) (1.540,-0.226) (1.580,-0.180) (1.610,-0.140) (1.637,-0.100) (1.660,-0.060) (1.680,-0.020) (1.697,0.020) (1.712,0.060) (1.724,0.100) (1.734,0.140) (1.742,0.180) (1.747,0.220) (1.750,0.260) (1.751,0.300) (1.750,0.340) (1.747,0.380) (1.741,0.420) (1.733,0.460) (1.722,0.500) (1.708,0.540) (1.692,0.580) (1.673,0.620) (1.650,0.660) (1.623,0.700) (1.593,0.740) (1.560,0.777) (1.520,0.816) (1.480,0.850) (1.441,0.880) (1.400,0.908) (1.360,0.932) (1.320,0.953) (1.280,0.973) (1.240,0.991) (1.200,1.007) (1.160,1.022) (1.120,1.036) (1.080,1.049) (1.040,1.061) (1.000,1.072) (0.960,1.082) (0.920,1.091) (0.880,1.099) (0.840,1.105) (0.800,1.109) (0.760,1.111) (0.720,1.111) (0.680,1.108) (0.640,1.102) (0.600,1.094) (0.560,1.083) (0.520,1.068) (0.480,1.051) (0.440,1.030) (0.400,1.005) (0.365,0.980) (0.320,0.942) (0.280,0.902) (0.244,0.860) (0.220,0.827) (0.190,0.780) (0.168,0.740) (0.150,0.700) (0.135,0.660) (0.123,0.620) (0.113,0.580) (0.107,0.540) (0.103,0.500) (0.102,0.460) (0.104,0.420) (0.108,0.380) (0.115,0.340) (0.126,0.300) (0.139,0.260) (0.156,0.220) (0.177,0.180) (0.200,0.141) (0.228,0.100) (0.259,0.060) (0.293,0.020) (0.320,-0.012) (0.360,-0.059) (0.394,-0.100) (0.420,-0.133) (0.458,-0.180) (0.492,-0.220) (0.520,-0.251) (0.550,-0.280) (0.580,-0.307) (0.620,-0.339) (0.660,-0.366) (0.700,-0.389) (0.740,-0.409) (0.780,-0.425) (0.820,-0.439) (0.860,-0.449) (0.900,-0.457) (0.940,-0.462)}
\EQ{(0.900,-0.320) (0.940,-0.322) (0.980,-0.321) (1.020,-0.317) (1.060,-0.310) (1.100,-0.301) (1.140,-0.288) (1.180,-0.273) (1.220,-0.254) (1.260,-0.232) (1.300,-0.205) (1.333,-0.180) (1.377,-0.140) (1.415,-0.100) (1.440,-0.069) (1.474,-0.020) (1.497,0.020) (1.516,0.060) (1.533,0.100) (1.546,0.140) (1.556,0.180) (1.564,0.220) (1.569,0.260) (1.571,0.300) (1.571,0.340) (1.568,0.380) (1.563,0.420) (1.555,0.460) (1.544,0.500) (1.530,0.540) (1.512,0.580) (1.491,0.620) (1.467,0.660) (1.440,0.696) (1.403,0.740) (1.362,0.780) (1.320,0.816) (1.288,0.840) (1.240,0.871) (1.200,0.894) (1.160,0.913) (1.120,0.931) (1.080,0.946) (1.040,0.959) (1.000,0.971) (0.961,0.980) (0.920,0.988) (0.880,0.993) (0.840,0.996) (0.800,0.997) (0.760,0.996) (0.720,0.992) (0.680,0.985) (0.640,0.976) (0.600,0.963) (0.560,0.947) (0.520,0.927) (0.480,0.903) (0.440,0.874) (0.400,0.840) (0.361,0.800) (0.328,0.760) (0.301,0.720) (0.280,0.684) (0.260,0.643) (0.243,0.600) (0.231,0.560) (0.222,0.520) (0.216,0.480) (0.214,0.440) (0.214,0.400) (0.217,0.360) (0.223,0.320) (0.232,0.280) (0.244,0.240) (0.260,0.200) (0.279,0.160) (0.300,0.122) (0.327,0.080) (0.356,0.040) (0.387,-0.000) (0.420,-0.039) (0.456,-0.080) (0.494,-0.120) (0.536,-0.160) (0.580,-0.196) (0.614,-0.220) (0.660,-0.248) (0.700,-0.268) (0.740,-0.285) (0.780,-0.298) (0.820,-0.309) (0.860,-0.316) (0.900,-0.320)}

    % Disco cargado, resto neutros - líneas de campo
% Generado por scripts/generar_datos.py (no editar a mano).
\LC{(-4.174,-3.059) (-4.201,-3.082) (-4.231,-3.103) (-4.262,-3.121) (-4.294,-3.138) (-4.326,-3.153) (-4.360,-3.167) (-4.393,-3.180) (-4.427,-3.192) (-4.461,-3.204) (-4.496,-3.214) (-4.530,-3.224) (-4.565,-3.233) (-4.600,-3.242) (-4.635,-3.249) (-4.670,-3.257) (-4.706,-3.263) (-4.741,-3.269) (-4.777,-3.274) (-4.813,-3.278) (-4.849,-3.282) (-4.884,-3.285) (-4.920,-3.287) (-4.950,-3.288)}
\LC{(-3.897,-3.290) (-3.920,-3.317) (-3.944,-3.344) (-3.969,-3.371) (-3.993,-3.397) (-4.018,-3.423) (-4.044,-3.448) (-4.069,-3.474) (-4.095,-3.499) (-4.121,-3.523) (-4.147,-3.548) (-4.174,-3.572) (-4.201,-3.597) (-4.227,-3.621) (-4.254,-3.644) (-4.282,-3.668) (-4.309,-3.691) (-4.336,-3.715) (-4.364,-3.738) (-4.392,-3.760) (-4.420,-3.783) (-4.448,-3.805) (-4.477,-3.827) (-4.505,-3.849) (-4.534,-3.871) (-4.564,-3.892) (-4.593,-3.912) (-4.623,-3.932) (-4.653,-3.952) (-4.684,-3.970) (-4.715,-3.988) (-4.747,-4.005) (-4.780,-4.021) (-4.813,-4.035) (-4.847,-4.047) (-4.881,-4.058) (-4.916,-4.066) (-4.950,-4.071)}
\LC{(-3.565,-3.567) (-3.587,-3.595) (-3.609,-3.624) (-3.630,-3.653) (-3.651,-3.682) (-3.672,-3.711) (-3.692,-3.741) (-3.712,-3.771) (-3.731,-3.801) (-3.750,-3.832) (-3.768,-3.863) (-3.786,-3.894) (-3.803,-3.926) (-3.820,-3.958) (-3.835,-3.991) (-3.850,-4.023) (-3.865,-4.056) (-3.878,-4.090) (-3.891,-4.123) (-3.903,-4.157) (-3.914,-4.192) (-3.924,-4.226) (-3.933,-4.261) (-3.940,-4.296) (-3.947,-4.332) (-3.952,-4.367) (-3.956,-4.403) (-3.959,-4.439) (-3.960,-4.450)}
\LC{(-3.275,-3.808) (-3.297,-3.837) (-3.316,-3.867) (-3.334,-3.898) (-3.351,-3.930) (-3.367,-3.962) (-3.381,-3.995) (-3.394,-4.029) (-3.407,-4.063) (-3.418,-4.097) (-3.428,-4.132) (-3.437,-4.166) (-3.445,-4.201) (-3.452,-4.237) (-3.458,-4.272) (-3.463,-4.308) (-3.468,-4.344) (-3.471,-4.379) (-3.474,-4.415) (-3.475,-4.450)}
\LC{(-3.085,-3.967) (-3.091,-4.002) (-3.096,-4.038) (-3.099,-4.074) (-3.101,-4.109) (-3.103,-4.145) (-3.105,-4.181) (-3.106,-4.217) (-3.107,-4.253) (-3.108,-4.289) (-3.109,-4.325) (-3.109,-4.361) (-3.110,-4.397) (-3.110,-4.433) (-3.110,-4.450)}
\LC{(-2.877,-3.803) (-2.851,-3.829) (-2.828,-3.857) (-2.807,-3.886) (-2.788,-3.916) (-2.770,-3.948) (-2.754,-3.980) (-2.739,-4.013) (-2.726,-4.046) (-2.714,-4.080) (-2.703,-4.114) (-2.693,-4.149) (-2.685,-4.184) (-2.677,-4.219) (-2.671,-4.254) (-2.665,-4.290) (-2.660,-4.326) (-2.657,-4.362) (-2.654,-4.397) (-2.652,-4.433) (-2.651,-4.450)}
\LC{(-2.604,-3.489) (-2.577,-3.513) (-2.552,-3.539) (-2.528,-3.565) (-2.504,-3.593) (-2.481,-3.621) (-2.459,-3.649) (-2.438,-3.678) (-2.418,-3.708) (-2.399,-3.739) (-2.381,-3.770) (-2.363,-3.801) (-2.347,-3.833) (-2.331,-3.866) (-2.316,-3.898) (-2.302,-3.932) (-2.289,-3.965) (-2.277,-3.999) (-2.266,-4.033) (-2.256,-4.068) (-2.246,-4.102) (-2.237,-4.137) (-2.229,-4.172) (-2.222,-4.208) (-2.216,-4.243) (-2.211,-4.279) (-2.206,-4.314) (-2.202,-4.350) (-2.199,-4.386) (-2.197,-4.422) (-2.196,-4.450)}
\LC{(-2.163,-2.981) (-2.137,-3.006) (-2.112,-3.031) (-2.086,-3.057) (-2.062,-3.083) (-2.038,-3.110) (-2.014,-3.137) (-1.992,-3.165) (-1.969,-3.194) (-1.948,-3.222) (-1.927,-3.252) (-1.906,-3.281) (-1.887,-3.311) (-1.867,-3.342) (-1.849,-3.373) (-1.831,-3.404) (-1.814,-3.435) (-1.797,-3.467) (-1.781,-3.500) (-1.765,-3.532) (-1.750,-3.565) (-1.736,-3.598) (-1.722,-3.631) (-1.709,-3.665) (-1.697,-3.698) (-1.685,-3.732) (-1.673,-3.766) (-1.662,-3.801) (-1.652,-3.835) (-1.642,-3.870) (-1.633,-3.905) (-1.624,-3.940) (-1.616,-3.975) (-1.609,-4.010) (-1.602,-4.045) (-1.595,-4.081) (-1.589,-4.116) (-1.584,-4.152) (-1.579,-4.187) (-1.575,-4.223) (-1.571,-4.259) (-1.568,-4.295) (-1.565,-4.331) (-1.563,-4.367) (-1.561,-4.403) (-1.560,-4.439) (-1.560,-4.450)}
\LC{(-1.349,-2.043) (-1.324,-2.068) (-1.300,-2.095) (-1.277,-2.122) (-1.255,-2.151) (-1.234,-2.180) (-1.214,-2.210) (-1.195,-2.241) (-1.177,-2.272) (-1.159,-2.303) (-1.143,-2.335) (-1.127,-2.368) (-1.112,-2.400) (-1.097,-2.433) (-1.083,-2.466) (-1.069,-2.500) (-1.056,-2.533) (-1.043,-2.567) (-1.031,-2.601) (-1.019,-2.635) (-1.007,-2.669) (-0.996,-2.703) (-0.985,-2.737) (-0.974,-2.771) (-0.964,-2.806) (-0.953,-2.840) (-0.943,-2.875) (-0.933,-2.910) (-0.924,-2.944) (-0.915,-2.979) (-0.905,-3.014) (-0.897,-3.049) (-0.888,-3.084) (-0.879,-3.119) (-0.871,-3.154) (-0.863,-3.189) (-0.855,-3.224) (-0.847,-3.259) (-0.840,-3.294) (-0.833,-3.330) (-0.826,-3.365) (-0.819,-3.400) (-0.812,-3.436) (-0.806,-3.471) (-0.799,-3.506) (-0.793,-3.542) (-0.787,-3.577) (-0.782,-3.613) (-0.776,-3.649) (-0.771,-3.684) (-0.766,-3.720) (-0.761,-3.756) (-0.757,-3.791) (-0.752,-3.827) (-0.748,-3.863) (-0.744,-3.898) (-0.740,-3.934) (-0.737,-3.970) (-0.734,-4.006) (-0.731,-4.042) (-0.728,-4.078) (-0.725,-4.114) (-0.723,-4.150) (-0.721,-4.186) (-0.719,-4.221) (-0.717,-4.257) (-0.716,-4.293) (-0.714,-4.329) (-0.713,-4.365) (-0.713,-4.401) (-0.712,-4.437) (-0.712,-4.450)}
\LC{(-2.983,-2.123) (-3.003,-2.094) (-3.024,-2.064) (-3.045,-2.035) (-3.067,-2.007) (-3.090,-1.979) (-3.113,-1.951) (-3.136,-1.924) (-3.161,-1.897) (-3.185,-1.871) (-3.211,-1.846) (-3.236,-1.820) (-3.263,-1.796) (-3.289,-1.772) (-3.317,-1.748) (-3.344,-1.725) (-3.372,-1.703) (-3.401,-1.681) (-3.430,-1.659) (-3.459,-1.638) (-3.489,-1.618) (-3.518,-1.598) (-3.549,-1.578) (-3.579,-1.559) (-3.610,-1.541) (-3.642,-1.523) (-3.673,-1.506) (-3.705,-1.489) (-3.737,-1.473) (-3.769,-1.457) (-3.802,-1.441) (-3.835,-1.426) (-3.868,-1.412) (-3.901,-1.398) (-3.934,-1.385) (-3.968,-1.372) (-4.002,-1.359) (-4.036,-1.347) (-4.070,-1.336) (-4.104,-1.325) (-4.138,-1.314) (-4.173,-1.304) (-4.208,-1.294) (-4.242,-1.285) (-4.277,-1.276) (-4.312,-1.268) (-4.347,-1.260) (-4.383,-1.252) (-4.418,-1.245) (-4.453,-1.238) (-4.489,-1.232) (-4.524,-1.226) (-4.560,-1.221) (-4.595,-1.216) (-4.631,-1.212) (-4.667,-1.208) (-4.703,-1.204) (-4.739,-1.201) (-4.775,-1.198) (-4.810,-1.196) (-4.846,-1.194) (-4.882,-1.192) (-4.918,-1.191) (-4.950,-1.190)}
\LC{(-3.751,-2.600) (-3.772,-2.570) (-3.793,-2.541) (-3.816,-2.513) (-3.839,-2.486) (-3.864,-2.459) (-3.889,-2.434) (-3.915,-2.409) (-3.942,-2.385) (-3.969,-2.361) (-3.997,-2.339) (-4.026,-2.317) (-4.055,-2.297) (-4.085,-2.277) (-4.116,-2.258) (-4.147,-2.239) (-4.178,-2.222) (-4.210,-2.205) (-4.243,-2.190) (-4.276,-2.175) (-4.309,-2.161) (-4.342,-2.147) (-4.376,-2.135) (-4.410,-2.123) (-4.444,-2.112) (-4.479,-2.101) (-4.513,-2.092) (-4.548,-2.083) (-4.583,-2.075) (-4.618,-2.067) (-4.654,-2.061) (-4.689,-2.055) (-4.725,-2.049) (-4.761,-2.044) (-4.796,-2.040) (-4.832,-2.037) (-4.868,-2.034) (-4.904,-2.032) (-4.940,-2.031) (-4.950,-2.030)}
\LC{(-4.160,-2.853) (-4.184,-2.827) (-4.212,-2.803) (-4.241,-2.783) (-4.272,-2.765) (-4.305,-2.749) (-4.338,-2.734) (-4.371,-2.721) (-4.405,-2.710) (-4.440,-2.700) (-4.475,-2.690) (-4.510,-2.682) (-4.545,-2.675) (-4.580,-2.668) (-4.616,-2.662) (-4.651,-2.657) (-4.687,-2.652) (-4.723,-2.648) (-4.759,-2.644) (-4.795,-2.641) (-4.830,-2.639) (-4.866,-2.637) (-4.902,-2.636) (-4.938,-2.635) (-4.950,-2.634)}
\LC{(-0.858,2.885) (-0.858,2.921) (-0.857,2.957) (-0.857,2.993) (-0.856,3.029) (-0.855,3.065) (-0.855,3.101) (-0.854,3.137) (-0.854,3.173) (-0.854,3.209) (-0.854,3.245) (-0.854,3.281) (-0.854,3.317) (-0.854,3.353) (-0.855,3.389) (-0.855,3.425) (-0.856,3.461) (-0.856,3.497) (-0.857,3.533) (-0.858,3.569) (-0.858,3.605) (-0.859,3.641) (-0.860,3.677) (-0.861,3.713) (-0.861,3.749) (-0.862,3.785) (-0.863,3.821) (-0.863,3.857) (-0.864,3.893) (-0.865,3.929) (-0.865,3.965) (-0.866,4.001) (-0.867,4.037) (-0.867,4.073) (-0.868,4.109) (-0.868,4.145) (-0.869,4.181) (-0.869,4.217) (-0.869,4.253) (-0.870,4.289) (-0.870,4.325) (-0.870,4.361) (-0.870,4.397) (-0.871,4.433) (-0.871,4.450)}
\LC{(-1.716,2.885) (-1.716,2.921) (-1.717,2.957) (-1.717,2.993) (-1.718,3.029) (-1.719,3.065) (-1.720,3.101) (-1.721,3.137) (-1.722,3.173) (-1.723,3.209) (-1.725,3.245) (-1.726,3.281) (-1.728,3.317) (-1.729,3.353) (-1.731,3.389) (-1.733,3.425) (-1.735,3.461) (-1.737,3.497) (-1.739,3.533) (-1.741,3.568) (-1.743,3.604) (-1.745,3.640) (-1.746,3.676) (-1.748,3.712) (-1.750,3.748) (-1.752,3.784) (-1.754,3.820) (-1.756,3.856) (-1.758,3.892) (-1.760,3.928) (-1.761,3.964) (-1.763,4.000) (-1.764,4.036) (-1.766,4.072) (-1.767,4.108) (-1.768,4.144) (-1.769,4.180) (-1.770,4.216) (-1.771,4.252) (-1.772,4.288) (-1.773,4.324) (-1.773,4.360) (-1.774,4.396) (-1.774,4.432) (-1.774,4.450)}
\LC{(-2.485,2.885) (-2.488,2.921) (-2.491,2.957) (-2.496,2.992) (-2.502,3.028) (-2.508,3.063) (-2.515,3.099) (-2.523,3.134) (-2.532,3.169) (-2.541,3.204) (-2.550,3.238) (-2.560,3.273) (-2.570,3.308) (-2.580,3.342) (-2.590,3.377) (-2.600,3.411) (-2.610,3.446) (-2.621,3.480) (-2.631,3.515) (-2.641,3.549) (-2.651,3.584) (-2.660,3.619) (-2.670,3.653) (-2.679,3.688) (-2.688,3.723) (-2.697,3.758) (-2.706,3.793) (-2.714,3.828) (-2.722,3.863) (-2.730,3.898) (-2.737,3.933) (-2.744,3.969) (-2.751,4.004) (-2.757,4.040) (-2.763,4.075) (-2.768,4.111) (-2.773,4.146) (-2.778,4.182) (-2.782,4.218) (-2.785,4.254) (-2.788,4.290) (-2.791,4.325) (-2.793,4.361) (-2.795,4.397) (-2.796,4.433) (-2.797,4.450)}
\LC{(-2.785,2.594) (-2.821,2.597) (-2.857,2.601) (-2.892,2.607) (-2.927,2.615) (-2.962,2.623) (-2.997,2.633) (-3.031,2.644) (-3.065,2.655) (-3.099,2.668) (-3.133,2.681) (-3.166,2.695) (-3.199,2.709) (-3.232,2.724) (-3.265,2.739) (-3.297,2.754) (-3.329,2.770) (-3.361,2.787) (-3.393,2.803) (-3.425,2.820) (-3.457,2.837) (-3.488,2.855) (-3.520,2.872) (-3.551,2.890) (-3.583,2.907) (-3.614,2.925) (-3.645,2.943) (-3.676,2.961) (-3.707,2.980) (-3.738,2.998) (-3.769,3.016) (-3.800,3.034) (-3.831,3.053) (-3.862,3.071) (-3.893,3.089) (-3.924,3.108) (-3.955,3.126) (-3.986,3.144) (-4.018,3.162) (-4.049,3.180) (-4.080,3.198) (-4.111,3.215) (-4.143,3.233) (-4.174,3.250) (-4.206,3.267) (-4.238,3.284) (-4.270,3.301) (-4.302,3.317) (-4.334,3.333) (-4.366,3.349) (-4.399,3.364) (-4.432,3.379) (-4.465,3.393) (-4.498,3.407) (-4.532,3.421) (-4.565,3.433) (-4.599,3.445) (-4.633,3.457) (-4.668,3.467) (-4.703,3.477) (-4.737,3.485) (-4.773,3.493) (-4.808,3.500) (-4.843,3.506) (-4.879,3.510) (-4.915,3.514) (-4.950,3.516)}
\LC{(-2.785,1.698) (-2.821,1.698) (-2.857,1.698) (-2.893,1.697) (-2.929,1.697) (-2.965,1.697) (-3.001,1.697) (-3.037,1.697) (-3.073,1.697) (-3.109,1.697) (-3.145,1.697) (-3.181,1.697) (-3.217,1.697) (-3.253,1.698) (-3.289,1.699) (-3.325,1.700) (-3.361,1.701) (-3.397,1.702) (-3.433,1.703) (-3.469,1.704) (-3.505,1.706) (-3.541,1.708) (-3.577,1.710) (-3.613,1.712) (-3.649,1.714) (-3.685,1.716) (-3.720,1.718) (-3.756,1.721) (-3.792,1.723) (-3.828,1.726) (-3.864,1.728) (-3.900,1.731) (-3.936,1.734) (-3.972,1.736) (-4.008,1.739) (-4.044,1.742) (-4.080,1.745) (-4.115,1.747) (-4.151,1.750) (-4.187,1.753) (-4.223,1.755) (-4.259,1.758) (-4.295,1.760) (-4.331,1.763) (-4.367,1.765) (-4.403,1.767) (-4.439,1.769) (-4.475,1.772) (-4.511,1.773) (-4.547,1.775) (-4.582,1.777) (-4.618,1.779) (-4.654,1.780) (-4.690,1.781) (-4.726,1.783) (-4.762,1.784) (-4.798,1.785) (-4.834,1.785) (-4.870,1.786) (-4.906,1.787) (-4.942,1.787) (-4.950,1.787)}
\LC{(1.361,0.418) (1.396,0.423) (1.432,0.429) (1.467,0.435) (1.502,0.442) (1.538,0.450) (1.573,0.458) (1.608,0.466) (1.643,0.475) (1.677,0.485) (1.712,0.495) (1.746,0.505) (1.781,0.516) (1.815,0.527) (1.849,0.539) (1.883,0.551) (1.917,0.563) (1.950,0.576) (1.984,0.589) (2.017,0.603) (2.050,0.616) (2.084,0.630) (2.117,0.645) (2.150,0.659) (2.182,0.674) (2.215,0.689) (2.248,0.703) (2.281,0.718) (2.314,0.733) (2.347,0.748) (2.380,0.762) (2.413,0.776) (2.446,0.790) (2.480,0.803) (2.513,0.815) (2.548,0.826) (2.582,0.837) (2.617,0.845) (2.652,0.853) (2.688,0.858) (2.724,0.861) (2.760,0.862) (2.796,0.861) (2.832,0.857) (2.867,0.851) (2.902,0.843) (2.937,0.834) (2.971,0.823) (3.005,0.811) (3.039,0.798) (3.072,0.784) (3.105,0.770) (3.138,0.755) (3.171,0.740) (3.203,0.725) (3.236,0.710) (3.269,0.695) (3.301,0.680) (3.334,0.665) (3.367,0.650) (3.400,0.635) (3.433,0.621) (3.466,0.607) (3.499,0.593) (3.533,0.580) (3.566,0.567) (3.600,0.554) (3.634,0.542) (3.667,0.529) (3.702,0.518) (3.736,0.506) (3.770,0.495) (3.804,0.485) (3.839,0.474) (3.873,0.464) (3.908,0.455) (3.943,0.446) (3.978,0.437) (4.013,0.428) (4.048,0.420) (4.083,0.412) (4.118,0.404) (4.153,0.397) (4.189,0.390) (4.224,0.384) (4.260,0.377) (4.295,0.371) (4.331,0.366) (4.366,0.360) (4.402,0.355) (4.438,0.351) (4.473,0.346) (4.509,0.342) (4.545,0.338) (4.581,0.335) (4.616,0.331) (4.652,0.328) (4.688,0.326) (4.724,0.323) (4.760,0.321) (4.796,0.320) (4.832,0.318) (4.868,0.317) (4.904,0.316) (4.940,0.315) (4.950,0.315)}
\LC{(1.331,0.535) (1.364,0.549) (1.397,0.564) (1.429,0.579) (1.461,0.596) (1.493,0.613) (1.524,0.630) (1.555,0.649) (1.586,0.668) (1.615,0.688) (1.645,0.709) (1.674,0.731) (1.702,0.753) (1.729,0.776) (1.756,0.800) (1.782,0.825) (1.807,0.851) (1.832,0.877) (1.855,0.905) (1.877,0.933) (1.898,0.962) (1.918,0.992) (1.937,1.023) (1.954,1.055) (1.969,1.087) (1.982,1.121) (1.994,1.155) (2.004,1.190) (2.011,1.225) (2.016,1.260) (2.019,1.296) (2.020,1.320)}
\LC{(1.274,0.642) (1.303,0.664) (1.331,0.687) (1.358,0.710) (1.385,0.734) (1.411,0.759) (1.436,0.785) (1.460,0.811) (1.484,0.838) (1.507,0.867) (1.528,0.895) (1.549,0.925) (1.568,0.955) (1.586,0.987) (1.602,1.019) (1.617,1.051) (1.631,1.085) (1.643,1.119) (1.653,1.153) (1.661,1.188) (1.667,1.224) (1.672,1.259) (1.674,1.295) (1.675,1.331)}
\LC{(1.204,0.724) (1.228,0.751) (1.251,0.779) (1.273,0.807) (1.295,0.836) (1.315,0.865) (1.335,0.896) (1.353,0.926) (1.371,0.958) (1.387,0.990) (1.402,1.023) (1.416,1.056) (1.428,1.090) (1.438,1.124) (1.447,1.159) (1.455,1.194) (1.460,1.230) (1.464,1.266) (1.466,1.302) (1.466,1.326)}
\LC{(1.130,0.782) (1.149,0.813) (1.167,0.844) (1.185,0.875) (1.202,0.907) (1.218,0.939) (1.233,0.972) (1.247,1.005) (1.260,1.039) (1.272,1.072) (1.283,1.107) (1.292,1.142) (1.300,1.177) (1.307,1.212) (1.311,1.248) (1.314,1.284) (1.315,1.320) (1.315,1.332)}
\LC{(1.047,0.826) (1.060,0.860) (1.073,0.893) (1.085,0.927) (1.098,0.961) (1.110,0.994) (1.122,1.029) (1.133,1.063) (1.143,1.097) (1.153,1.132) (1.162,1.167) (1.170,1.202) (1.177,1.237) (1.181,1.273) (1.183,1.309) (1.183,1.321)}
\LC{(0.957,0.854) (0.965,0.889) (0.973,0.924) (0.981,0.959) (0.990,0.994) (0.998,1.029) (1.008,1.064) (1.017,1.098) (1.027,1.133) (1.038,1.167) (1.049,1.202) (1.061,1.236) (1.072,1.270) (1.081,1.305) (1.082,1.329)}
\LC{(0.863,0.865) (0.864,0.901) (0.867,0.937) (0.869,0.973) (0.873,1.008) (0.877,1.044) (0.882,1.080) (0.889,1.115) (0.896,1.150) (0.906,1.185) (0.917,1.219) (0.931,1.252) (0.948,1.284) (0.969,1.313) (0.995,1.338) (1.026,1.356) (1.062,1.363)}
\LC{(0.769,0.859) (0.764,0.894) (0.760,0.930) (0.756,0.966) (0.753,1.002) (0.751,1.038) (0.750,1.074) (0.749,1.110) (0.750,1.146) (0.752,1.182) (0.755,1.217) (0.760,1.253) (0.766,1.289) (0.773,1.324) (0.783,1.359) (0.794,1.393) (0.806,1.427) (0.820,1.460) (0.836,1.492) (0.853,1.524) (0.871,1.555) (0.890,1.586) (0.908,1.617) (0.926,1.648) (0.941,1.681) (0.952,1.715) (0.960,1.750) (0.965,1.786) (0.969,1.821) (0.972,1.857) (0.973,1.893) (0.974,1.929) (0.974,1.965) (0.974,2.001) (0.973,2.037) (0.972,2.073) (0.970,2.109) (0.968,2.145) (0.966,2.181) (0.963,2.217) (0.960,2.253) (0.957,2.289) (0.954,2.325) (0.951,2.360) (0.947,2.396) (0.944,2.432) (0.940,2.468) (0.936,2.504) (0.933,2.540) (0.929,2.575) (0.925,2.611) (0.921,2.647) (0.917,2.683) (0.914,2.719) (0.910,2.754) (0.906,2.790) (0.902,2.826) (0.899,2.862) (0.895,2.898) (0.892,2.933) (0.888,2.969) (0.885,3.005) (0.882,3.041) (0.879,3.077) (0.876,3.113) (0.873,3.149) (0.870,3.184) (0.867,3.220) (0.865,3.256) (0.862,3.292) (0.860,3.328) (0.857,3.364) (0.855,3.400) (0.853,3.436) (0.851,3.472) (0.849,3.508) (0.847,3.544) (0.845,3.580) (0.844,3.616) (0.842,3.652) (0.841,3.688) (0.839,3.724) (0.838,3.760) (0.837,3.795) (0.835,3.831) (0.834,3.867) (0.833,3.903) (0.832,3.939) (0.831,3.975) (0.831,4.011) (0.830,4.047) (0.829,4.083) (0.829,4.119) (0.828,4.155) (0.828,4.191) (0.827,4.227) (0.827,4.263) (0.826,4.299) (0.826,4.335) (0.826,4.371) (0.826,4.407) (0.826,4.443) (0.826,4.450)}
\LC{(0.665,0.831) (0.653,0.864) (0.641,0.898) (0.629,0.932) (0.618,0.967) (0.607,1.001) (0.596,1.035) (0.586,1.070) (0.576,1.105) (0.567,1.139) (0.558,1.174) (0.550,1.209) (0.541,1.244) (0.533,1.279) (0.526,1.314) (0.518,1.350) (0.510,1.385) (0.503,1.420) (0.495,1.455) (0.487,1.490) (0.479,1.525) (0.471,1.560) (0.462,1.595) (0.453,1.630) (0.443,1.665) (0.433,1.699) (0.422,1.734) (0.411,1.768) (0.399,1.802) (0.387,1.836) (0.375,1.870) (0.361,1.903) (0.348,1.936) (0.334,1.970) (0.320,2.003) (0.305,2.035) (0.290,2.068) (0.274,2.101) (0.259,2.133) (0.243,2.165) (0.227,2.198) (0.211,2.230) (0.194,2.262) (0.178,2.294) (0.161,2.326) (0.144,2.358) (0.128,2.390) (0.111,2.422) (0.095,2.454) (0.079,2.486) (0.063,2.518) (0.047,2.551) (0.032,2.583) (0.017,2.616) (0.003,2.649) (-0.011,2.682) (-0.024,2.716) (-0.037,2.749) (-0.049,2.784) (-0.060,2.818) (-0.070,2.852) (-0.079,2.887) (-0.088,2.922) (-0.096,2.957) (-0.103,2.992) (-0.110,3.028) (-0.116,3.063) (-0.122,3.099) (-0.127,3.134) (-0.132,3.170) (-0.137,3.206) (-0.141,3.241) (-0.145,3.277) (-0.148,3.313) (-0.152,3.349) (-0.155,3.385) (-0.158,3.421) (-0.161,3.456) (-0.163,3.492) (-0.166,3.528) (-0.168,3.564) (-0.170,3.600) (-0.172,3.636) (-0.174,3.672) (-0.176,3.708) (-0.178,3.744) (-0.180,3.780) (-0.181,3.816) (-0.183,3.852) (-0.184,3.888) (-0.185,3.924) (-0.187,3.960) (-0.188,3.996) (-0.189,4.032) (-0.190,4.068) (-0.190,4.104) (-0.191,4.140) (-0.192,4.176) (-0.193,4.212) (-0.193,4.248) (-0.194,4.284) (-0.194,4.320) (-0.194,4.356) (-0.195,4.392) (-0.195,4.428) (-0.195,4.450)}
\LC{(0.581,0.789) (0.561,0.819) (0.543,0.850) (0.524,0.881) (0.506,0.912) (0.487,0.943) (0.469,0.974) (0.451,1.005) (0.433,1.036) (0.415,1.068) (0.397,1.099) (0.379,1.130) (0.361,1.161) (0.343,1.192) (0.325,1.223) (0.307,1.254) (0.288,1.285) (0.269,1.316) (0.250,1.346) (0.230,1.376) (0.211,1.406) (0.190,1.436) (0.169,1.466) (0.148,1.495) (0.126,1.523) (0.104,1.552) (0.081,1.579) (0.058,1.607) (0.034,1.634) (0.009,1.660) (-0.016,1.686) (-0.041,1.711) (-0.068,1.735) (-0.095,1.759) (-0.122,1.782) (-0.151,1.805) (-0.180,1.826) (-0.209,1.847) (-0.239,1.866) (-0.270,1.885) (-0.302,1.902) (-0.334,1.919) (-0.366,1.934) (-0.400,1.948) (-0.433,1.960) (-0.468,1.971) (-0.503,1.980) (-0.538,1.987) (-0.573,1.993) (-0.609,1.996) (-0.645,1.998) (-0.657,1.998)}
\LC{(0.495,0.723) (0.470,0.749) (0.445,0.775) (0.420,0.801) (0.395,0.827) (0.370,0.852) (0.345,0.878) (0.319,0.904) (0.294,0.929) (0.268,0.955) (0.243,0.980) (0.217,1.005) (0.191,1.030) (0.165,1.055) (0.139,1.080) (0.113,1.105) (0.086,1.129) (0.060,1.153) (0.033,1.177) (0.005,1.200) (-0.022,1.224) (-0.050,1.247) (-0.078,1.269) (-0.106,1.291) (-0.135,1.313) (-0.165,1.333) (-0.194,1.354) (-0.224,1.374) (-0.255,1.392) (-0.286,1.411) (-0.318,1.428) (-0.350,1.444) (-0.383,1.459) (-0.416,1.472) (-0.450,1.485) (-0.484,1.495) (-0.519,1.504) (-0.554,1.511) (-0.590,1.516) (-0.626,1.519) (-0.662,1.520)}
\LC{(0.425,0.641) (0.396,0.662) (0.366,0.682) (0.336,0.702) (0.305,0.721) (0.275,0.740) (0.244,0.759) (0.214,0.778) (0.183,0.797) (0.152,0.815) (0.121,0.833) (0.090,0.851) (0.058,0.869) (0.027,0.887) (-0.004,0.905) (-0.036,0.922) (-0.067,0.940) (-0.098,0.957) (-0.130,0.975) (-0.161,0.993) (-0.193,1.010) (-0.224,1.028) (-0.256,1.045) (-0.287,1.063) (-0.318,1.081) (-0.349,1.099) (-0.380,1.117) (-0.411,1.135) (-0.442,1.154) (-0.473,1.173) (-0.504,1.191) (-0.534,1.210) (-0.565,1.229) (-0.596,1.248) (-0.628,1.264) (-0.651,1.272)}
\LC{(0.374,0.547) (0.340,0.559) (0.306,0.572) (0.273,0.584) (0.238,0.595) (0.204,0.606) (0.170,0.617) (0.135,0.627) (0.101,0.637) (0.066,0.646) (0.031,0.655) (-0.004,0.664) (-0.039,0.673) (-0.074,0.682) (-0.109,0.690) (-0.144,0.699) (-0.179,0.707) (-0.214,0.716) (-0.248,0.725) (-0.283,0.734) (-0.318,0.744) (-0.352,0.754) (-0.387,0.764) (-0.421,0.776) (-0.455,0.788) (-0.489,0.801) (-0.522,0.814) (-0.555,0.829) (-0.587,0.845) (-0.618,0.862) (-0.649,0.881) (-0.679,0.901) (-0.708,0.923) (-0.735,0.946) (-0.762,0.971) (-0.786,0.997) (-0.808,1.026) (-0.828,1.055) (-0.846,1.087) (-0.861,1.119) (-0.873,1.153) (-0.882,1.188) (-0.887,1.224) (-0.889,1.260)}
\LC{(0.346,0.457) (0.311,0.463) (0.275,0.469) (0.240,0.473) (0.204,0.477) (0.168,0.480) (0.132,0.482) (0.096,0.484) (0.060,0.485) (0.024,0.485) (-0.012,0.484) (-0.048,0.483) (-0.084,0.481) (-0.120,0.479) (-0.156,0.477) (-0.192,0.474) (-0.228,0.470) (-0.263,0.467) (-0.299,0.463) (-0.335,0.460) (-0.371,0.456) (-0.407,0.452) (-0.442,0.448) (-0.478,0.444) (-0.514,0.441) (-0.550,0.438) (-0.586,0.435) (-0.622,0.432) (-0.658,0.430) (-0.694,0.428) (-0.730,0.427) (-0.766,0.426) (-0.802,0.426) (-0.838,0.427) (-0.874,0.428) (-0.909,0.430) (-0.945,0.432) (-0.981,0.436) (-1.017,0.440) (-1.053,0.445) (-1.088,0.451) (-1.123,0.458) (-1.159,0.465) (-1.194,0.473) (-1.229,0.483) (-1.263,0.493) (-1.297,0.503) (-1.332,0.515) (-1.365,0.527) (-1.399,0.540) (-1.432,0.554) (-1.465,0.569) (-1.498,0.584) (-1.530,0.599) (-1.562,0.616) (-1.594,0.632) (-1.626,0.650) (-1.657,0.667) (-1.688,0.685) (-1.719,0.704) (-1.750,0.722) (-1.781,0.741) (-1.811,0.760) (-1.842,0.779) (-1.872,0.798) (-1.903,0.817) (-1.934,0.835) (-1.965,0.853) (-1.997,0.871) (-2.029,0.887) (-2.062,0.902) (-2.095,0.915) (-2.129,0.926) (-2.164,0.935) (-2.200,0.940) (-2.236,0.942) (-2.272,0.940) (-2.308,0.936) (-2.343,0.929) (-2.378,0.920) (-2.412,0.910) (-2.447,0.899) (-2.481,0.887) (-2.515,0.875) (-2.548,0.863) (-2.582,0.851) (-2.616,0.839) (-2.651,0.828) (-2.685,0.817) (-2.719,0.807) (-2.754,0.797) (-2.789,0.788) (-2.824,0.780) (-2.859,0.772) (-2.894,0.764) (-2.930,0.757) (-2.965,0.751) (-3.001,0.745) (-3.036,0.740) (-3.072,0.735) (-3.108,0.731) (-3.143,0.727) (-3.179,0.723) (-3.215,0.720) (-3.251,0.717) (-3.287,0.715) (-3.323,0.712) (-3.359,0.710) (-3.395,0.709) (-3.431,0.707) (-3.467,0.706) (-3.503,0.705) (-3.539,0.704) (-3.575,0.704) (-3.611,0.703) (-3.647,0.703) (-3.683,0.703) (-3.719,0.703) (-3.755,0.703) (-3.791,0.704) (-3.827,0.704) (-3.863,0.704) (-3.899,0.705) (-3.935,0.706) (-3.971,0.706) (-4.007,0.707) (-4.043,0.708) (-4.079,0.708) (-4.115,0.709) (-4.151,0.710) (-4.186,0.711) (-4.222,0.712) (-4.258,0.713) (-4.294,0.714) (-4.330,0.714) (-4.366,0.715) (-4.402,0.716) (-4.438,0.717) (-4.474,0.718) (-4.510,0.718) (-4.546,0.719) (-4.582,0.720) (-4.618,0.720) (-4.654,0.721) (-4.690,0.721) (-4.726,0.722) (-4.762,0.722) (-4.798,0.723) (-4.834,0.723) (-4.870,0.723) (-4.906,0.723) (-4.942,0.724) (-4.950,0.724)}
\LC{(0.335,0.363) (0.299,0.362) (0.263,0.359) (0.227,0.356) (0.192,0.351) (0.156,0.345) (0.121,0.338) (0.086,0.330) (0.051,0.320) (0.017,0.309) (-0.017,0.297) (-0.050,0.284) (-0.083,0.269) (-0.116,0.253) (-0.147,0.236) (-0.178,0.218) (-0.209,0.199) (-0.239,0.179) (-0.267,0.157) (-0.295,0.135) (-0.323,0.111) (-0.349,0.086) (-0.374,0.060) (-0.398,0.034) (-0.421,0.006) (-0.443,-0.023) (-0.463,-0.052) (-0.482,-0.083) (-0.500,-0.114) (-0.516,-0.146) (-0.530,-0.179) (-0.543,-0.213) (-0.554,-0.247) (-0.563,-0.282) (-0.569,-0.318) (-0.574,-0.353) (-0.576,-0.389) (-0.576,-0.425) (-0.573,-0.461) (-0.567,-0.497) (-0.559,-0.532) (-0.547,-0.566) (-0.534,-0.599) (-0.518,-0.632) (-0.513,-0.642)}
\LC{(0.339,0.282) (0.304,0.276) (0.269,0.267) (0.235,0.257) (0.201,0.245) (0.167,0.232) (0.134,0.217) (0.102,0.200) (0.072,0.182) (0.042,0.162) (0.013,0.140) (-0.014,0.116) (-0.039,0.091) (-0.063,0.064) (-0.085,0.035) (-0.104,0.004) (-0.121,-0.027) (-0.135,-0.060) (-0.147,-0.094) (-0.155,-0.129) (-0.161,-0.165) (-0.163,-0.201) (-0.161,-0.237) (-0.156,-0.272) (-0.148,-0.307) (-0.136,-0.341) (-0.119,-0.373) (-0.105,-0.393)}
\LC{(0.356,0.203) (0.323,0.190) (0.290,0.176) (0.258,0.160) (0.226,0.142) (0.196,0.122) (0.168,0.100) (0.141,0.076) (0.116,0.050) (0.093,0.022) (0.074,-0.008) (0.057,-0.040) (0.045,-0.074) (0.036,-0.109) (0.033,-0.144) (0.034,-0.180) (0.041,-0.216) (0.052,-0.250) (0.067,-0.283)}
\LC{(0.385,0.128) (0.354,0.110) (0.323,0.091) (0.294,0.071) (0.265,0.049) (0.239,0.025) (0.214,-0.001) (0.191,-0.029) (0.172,-0.060) (0.158,-0.093) (0.149,-0.128) (0.149,-0.164) (0.157,-0.199) (0.165,-0.221)}
\LC{(0.418,0.069) (0.389,0.048) (0.360,0.026) (0.332,0.004) (0.305,-0.020) (0.279,-0.044) (0.254,-0.071) (0.232,-0.099) (0.215,-0.130) (0.205,-0.165) (0.206,-0.189)}
\LC{(0.450,0.026) (0.423,0.002) (0.395,-0.021) (0.367,-0.044) (0.340,-0.067) (0.312,-0.090) (0.285,-0.114) (0.259,-0.138) (0.233,-0.164) (0.209,-0.190)}
\LC{(0.496,-0.024) (0.470,-0.049) (0.443,-0.073) (0.416,-0.097) (0.388,-0.119) (0.359,-0.140) (0.329,-0.160) (0.297,-0.177) (0.263,-0.189) (0.228,-0.193) (0.206,-0.195)}
\LC{(0.548,-0.067) (0.524,-0.094) (0.499,-0.120) (0.473,-0.145) (0.446,-0.168) (0.417,-0.190) (0.387,-0.209) (0.354,-0.225) (0.321,-0.238) (0.285,-0.245) (0.250,-0.245) (0.215,-0.236) (0.184,-0.218)}
\LC{(0.617,-0.109) (0.597,-0.139) (0.576,-0.169) (0.554,-0.197) (0.530,-0.223) (0.504,-0.248) (0.476,-0.271) (0.447,-0.292) (0.415,-0.310) (0.383,-0.325) (0.349,-0.337) (0.314,-0.344) (0.278,-0.347) (0.242,-0.345) (0.206,-0.338) (0.173,-0.326) (0.141,-0.309) (0.113,-0.287)}
\LC{(0.679,-0.136) (0.664,-0.168) (0.647,-0.200) (0.629,-0.231) (0.609,-0.261) (0.587,-0.290) (0.563,-0.317) (0.538,-0.343) (0.511,-0.367) (0.483,-0.388) (0.452,-0.408) (0.421,-0.425) (0.388,-0.440) (0.354,-0.451) (0.319,-0.459) (0.283,-0.464) (0.247,-0.466) (0.211,-0.464) (0.176,-0.458) (0.141,-0.448) (0.108,-0.435) (0.076,-0.417) (0.047,-0.395) (0.038,-0.388)}
\LC{(0.770,-0.159) (0.762,-0.194) (0.752,-0.228) (0.741,-0.263) (0.728,-0.296) (0.714,-0.329) (0.698,-0.362) (0.680,-0.393) (0.661,-0.424) (0.641,-0.453) (0.619,-0.482) (0.596,-0.509) (0.571,-0.535) (0.545,-0.560) (0.517,-0.583) (0.489,-0.605) (0.459,-0.625) (0.428,-0.644) (0.396,-0.660) (0.363,-0.675) (0.329,-0.688) (0.295,-0.698) (0.260,-0.707) (0.224,-0.713) (0.189,-0.717) (0.153,-0.718) (0.117,-0.717) (0.081,-0.714) (0.045,-0.708) (0.011,-0.699) (-0.024,-0.688) (-0.057,-0.674) (-0.089,-0.657) (-0.120,-0.639) (-0.149,-0.618) (-0.159,-0.611)}
\LC{(0.864,-0.165) (0.863,-0.201) (0.860,-0.237) (0.857,-0.273) (0.852,-0.308) (0.845,-0.344) (0.838,-0.379) (0.829,-0.414) (0.820,-0.448) (0.809,-0.483) (0.797,-0.517) (0.783,-0.550) (0.769,-0.583) (0.753,-0.616) (0.737,-0.647) (0.719,-0.679) (0.700,-0.710) (0.681,-0.740) (0.660,-0.769) (0.639,-0.798) (0.617,-0.827) (0.594,-0.854) (0.570,-0.881) (0.545,-0.907) (0.519,-0.933) (0.493,-0.958) (0.467,-0.982) (0.439,-1.005) (0.411,-1.027) (0.382,-1.049) (0.353,-1.070) (0.323,-1.091) (0.293,-1.110) (0.262,-1.129) (0.231,-1.147) (0.200,-1.164) (0.168,-1.181) (0.135,-1.196) (0.103,-1.211) (0.069,-1.225) (0.036,-1.239) (0.002,-1.252) (-0.032,-1.263) (-0.066,-1.274) (-0.100,-1.285) (-0.135,-1.294) (-0.170,-1.302) (-0.205,-1.310) (-0.241,-1.317) (-0.276,-1.322) (-0.312,-1.327) (-0.348,-1.330) (-0.384,-1.332) (-0.420,-1.333) (-0.456,-1.332) (-0.492,-1.330) (-0.527,-1.326) (-0.563,-1.319) (-0.598,-1.310) (-0.631,-1.297) (-0.664,-1.282) (-0.694,-1.263) (-0.713,-1.248)}
\LC{(0.971,-0.151) (0.978,-0.186) (0.985,-0.221) (0.990,-0.257) (0.994,-0.293) (0.998,-0.328) (1.000,-0.364) (1.002,-0.400) (1.003,-0.436) (1.003,-0.472) (1.002,-0.508) (1.001,-0.544) (0.999,-0.580) (0.996,-0.616) (0.992,-0.652) (0.988,-0.688) (0.983,-0.723) (0.977,-0.759) (0.971,-0.794) (0.964,-0.830) (0.957,-0.865) (0.949,-0.900) (0.941,-0.935) (0.932,-0.970) (0.923,-1.005) (0.914,-1.040) (0.904,-1.074) (0.894,-1.109) (0.884,-1.143) (0.873,-1.178) (0.862,-1.212) (0.851,-1.246) (0.840,-1.280) (0.828,-1.315) (0.817,-1.349) (0.805,-1.383) (0.793,-1.417) (0.782,-1.451) (0.770,-1.485) (0.758,-1.519) (0.746,-1.553) (0.735,-1.587) (0.723,-1.621) (0.712,-1.655) (0.701,-1.690) (0.690,-1.724) (0.679,-1.758) (0.668,-1.793) (0.658,-1.827) (0.648,-1.862) (0.638,-1.896) (0.628,-1.931) (0.619,-1.966) (0.610,-2.001) (0.602,-2.036) (0.594,-2.071) (0.586,-2.106) (0.578,-2.141) (0.571,-2.176) (0.565,-2.212) (0.558,-2.247) (0.552,-2.283) (0.547,-2.318) (0.541,-2.354) (0.537,-2.389) (0.532,-2.425) (0.528,-2.461) (0.524,-2.497) (0.521,-2.533) (0.517,-2.568) (0.514,-2.604) (0.512,-2.640) (0.510,-2.676) (0.508,-2.712) (0.506,-2.748) (0.504,-2.784) (0.503,-2.820) (0.502,-2.856) (0.501,-2.892) (0.501,-2.928) (0.500,-2.964) (0.500,-3.000) (0.500,-3.036) (0.500,-3.072) (0.501,-3.108) (0.501,-3.144) (0.501,-3.180) (0.502,-3.216) (0.503,-3.252) (0.503,-3.288) (0.504,-3.324) (0.505,-3.360) (0.506,-3.396) (0.507,-3.432) (0.509,-3.468) (0.510,-3.504) (0.511,-3.540) (0.512,-3.576) (0.513,-3.612) (0.514,-3.648) (0.516,-3.684) (0.517,-3.720) (0.518,-3.756) (0.519,-3.792) (0.520,-3.828) (0.521,-3.864) (0.523,-3.900) (0.524,-3.936) (0.525,-3.972) (0.526,-4.008) (0.526,-4.044) (0.527,-4.080) (0.528,-4.116) (0.529,-4.152) (0.529,-4.188) (0.530,-4.224) (0.531,-4.260) (0.531,-4.296) (0.531,-4.332) (0.532,-4.368) (0.532,-4.404) (0.532,-4.440) (0.532,-4.450)}
\LC{(1.084,-0.109) (1.100,-0.141) (1.115,-0.174) (1.130,-0.207) (1.144,-0.240) (1.157,-0.273) (1.170,-0.307) (1.182,-0.341) (1.194,-0.375) (1.205,-0.409) (1.216,-0.443) (1.226,-0.478) (1.236,-0.512) (1.246,-0.547) (1.255,-0.582) (1.264,-0.617) (1.272,-0.652) (1.280,-0.687) (1.288,-0.722) (1.295,-0.757) (1.302,-0.793) (1.309,-0.828) (1.315,-0.863) (1.322,-0.899) (1.328,-0.934) (1.334,-0.970) (1.340,-1.005) (1.345,-1.041) (1.351,-1.076) (1.356,-1.112) (1.362,-1.148) (1.367,-1.183) (1.372,-1.219) (1.377,-1.255) (1.382,-1.290) (1.387,-1.326) (1.392,-1.362) (1.397,-1.397) (1.401,-1.433) (1.406,-1.469) (1.411,-1.504) (1.416,-1.540) (1.421,-1.576) (1.426,-1.611) (1.431,-1.647) (1.436,-1.682) (1.441,-1.718) (1.447,-1.754) (1.452,-1.789) (1.457,-1.825) (1.462,-1.861) (1.468,-1.896) (1.473,-1.932) (1.479,-1.967) (1.485,-2.003) (1.490,-2.038) (1.496,-2.074) (1.502,-2.109) (1.508,-2.145) (1.514,-2.180) (1.520,-2.216) (1.526,-2.251) (1.533,-2.287) (1.539,-2.322) (1.545,-2.358) (1.552,-2.393) (1.558,-2.429) (1.565,-2.464) (1.571,-2.499) (1.578,-2.535) (1.584,-2.570) (1.591,-2.605) (1.598,-2.641) (1.604,-2.676) (1.611,-2.712) (1.618,-2.747) (1.624,-2.782) (1.631,-2.818) (1.638,-2.853) (1.644,-2.888) (1.651,-2.924) (1.658,-2.959) (1.664,-2.995) (1.671,-3.030) (1.677,-3.065) (1.684,-3.101) (1.690,-3.136) (1.696,-3.172) (1.702,-3.207) (1.709,-3.243) (1.715,-3.278) (1.721,-3.314) (1.726,-3.349) (1.732,-3.385) (1.738,-3.420) (1.743,-3.456) (1.749,-3.491) (1.754,-3.527) (1.759,-3.563) (1.764,-3.598) (1.769,-3.634) (1.774,-3.670) (1.778,-3.705) (1.783,-3.741) (1.787,-3.777) (1.791,-3.813) (1.795,-3.848) (1.799,-3.884) (1.802,-3.920) (1.805,-3.956) (1.809,-3.992) (1.812,-4.028) (1.814,-4.064) (1.817,-4.100) (1.819,-4.135) (1.821,-4.171) (1.823,-4.207) (1.825,-4.243) (1.826,-4.279) (1.828,-4.315) (1.829,-4.351) (1.829,-4.387) (1.830,-4.423) (1.830,-4.450)}
\LC{(1.185,-0.041) (1.208,-0.069) (1.231,-0.096) (1.254,-0.125) (1.276,-0.153) (1.297,-0.182) (1.319,-0.210) (1.340,-0.240) (1.361,-0.269) (1.382,-0.298) (1.402,-0.328) (1.423,-0.358) (1.442,-0.388) (1.462,-0.418) (1.482,-0.448) (1.501,-0.478) (1.520,-0.509) (1.539,-0.540) (1.558,-0.570) (1.576,-0.601) (1.595,-0.632) (1.613,-0.663) (1.631,-0.694) (1.649,-0.725) (1.667,-0.757) (1.685,-0.788) (1.702,-0.819) (1.720,-0.851) (1.738,-0.882) (1.755,-0.914) (1.772,-0.945) (1.790,-0.977) (1.807,-1.008) (1.824,-1.040) (1.842,-1.072) (1.859,-1.103) (1.876,-1.135) (1.893,-1.166) (1.910,-1.198) (1.928,-1.230) (1.945,-1.261) (1.962,-1.293) (1.979,-1.325) (1.996,-1.356) (2.014,-1.388) (2.031,-1.419) (2.048,-1.451) (2.066,-1.482) (2.083,-1.514) (2.100,-1.545) (2.118,-1.577) (2.135,-1.608) (2.153,-1.640) (2.170,-1.671) (2.188,-1.703) (2.206,-1.734) (2.223,-1.765) (2.241,-1.797) (2.259,-1.828) (2.277,-1.859) (2.295,-1.890) (2.313,-1.922) (2.331,-1.953) (2.349,-1.984) (2.367,-2.015) (2.385,-2.046) (2.403,-2.077) (2.421,-2.108) (2.439,-2.140) (2.458,-2.171) (2.476,-2.202) (2.494,-2.233) (2.513,-2.264) (2.531,-2.295) (2.549,-2.325) (2.568,-2.356) (2.586,-2.387) (2.605,-2.418) (2.623,-2.449) (2.642,-2.480) (2.660,-2.511) (2.678,-2.542) (2.697,-2.573) (2.715,-2.604) (2.734,-2.635) (2.752,-2.665) (2.771,-2.696) (2.789,-2.727) (2.808,-2.758) (2.826,-2.789) (2.844,-2.820) (2.862,-2.851) (2.881,-2.882) (2.899,-2.913) (2.917,-2.945) (2.935,-2.976) (2.953,-3.007) (2.971,-3.038) (2.989,-3.069) (3.006,-3.101) (3.024,-3.132) (3.041,-3.164) (3.059,-3.195) (3.076,-3.227) (3.093,-3.258) (3.110,-3.290) (3.127,-3.322) (3.143,-3.354) (3.160,-3.386) (3.176,-3.418) (3.192,-3.450) (3.208,-3.483) (3.224,-3.515) (3.239,-3.548) (3.254,-3.580) (3.269,-3.613) (3.283,-3.646) (3.297,-3.679) (3.311,-3.713) (3.325,-3.746) (3.338,-3.779) (3.351,-3.813) (3.363,-3.847) (3.375,-3.881) (3.386,-3.915) (3.397,-3.949) (3.407,-3.984) (3.417,-4.019) (3.426,-4.053) (3.435,-4.088) (3.443,-4.123) (3.450,-4.159) (3.457,-4.194) (3.463,-4.229) (3.469,-4.265) (3.473,-4.301) (3.477,-4.337) (3.480,-4.372) (3.483,-4.408) (3.484,-4.444) (3.485,-4.450)}
\LC{(1.275,0.059) (1.304,0.038) (1.333,0.017) (1.362,-0.005) (1.391,-0.026) (1.420,-0.047) (1.449,-0.069) (1.478,-0.090) (1.507,-0.112) (1.535,-0.134) (1.564,-0.156) (1.592,-0.178) (1.620,-0.201) (1.648,-0.223) (1.676,-0.245) (1.704,-0.268) (1.732,-0.291) (1.760,-0.314) (1.788,-0.337) (1.815,-0.360) (1.843,-0.383) (1.870,-0.407) (1.898,-0.430) (1.925,-0.453) (1.952,-0.477) (1.979,-0.501) (2.006,-0.525) (2.033,-0.548) (2.060,-0.572) (2.087,-0.596) (2.114,-0.620) (2.141,-0.644) (2.167,-0.668) (2.194,-0.693) (2.221,-0.717) (2.247,-0.741) (2.274,-0.765) (2.301,-0.789) (2.327,-0.814) (2.354,-0.838) (2.380,-0.862) (2.407,-0.886) (2.434,-0.911) (2.460,-0.935) (2.487,-0.959) (2.514,-0.983) (2.540,-1.007) (2.567,-1.031) (2.594,-1.055) (2.621,-1.079) (2.648,-1.103) (2.675,-1.127) (2.702,-1.151) (2.729,-1.175) (2.756,-1.198) (2.783,-1.222) (2.811,-1.245) (2.838,-1.269) (2.865,-1.292) (2.893,-1.315) (2.920,-1.338) (2.948,-1.361) (2.976,-1.384) (3.004,-1.407) (3.032,-1.429) (3.060,-1.452) (3.088,-1.474) (3.116,-1.497) (3.145,-1.519) (3.173,-1.541) (3.202,-1.562) (3.231,-1.584) (3.260,-1.606) (3.289,-1.627) (3.318,-1.648) (3.347,-1.669) (3.377,-1.690) (3.406,-1.710) (3.436,-1.731) (3.466,-1.751) (3.495,-1.771) (3.526,-1.791) (3.556,-1.810) (3.586,-1.830) (3.617,-1.849) (3.647,-1.867) (3.678,-1.886) (3.709,-1.904) (3.740,-1.922) (3.772,-1.940) (3.803,-1.958) (3.835,-1.975) (3.867,-1.992) (3.899,-2.008) (3.931,-2.024) (3.963,-2.040) (3.996,-2.056) (4.028,-2.071) (4.061,-2.086) (4.094,-2.100) (4.127,-2.114) (4.160,-2.128) (4.194,-2.141) (4.228,-2.154) (4.261,-2.166) (4.295,-2.178) (4.330,-2.189) (4.364,-2.200) (4.398,-2.211) (4.433,-2.220) (4.468,-2.230) (4.503,-2.239) (4.538,-2.247) (4.573,-2.255) (4.608,-2.262) (4.644,-2.268) (4.679,-2.274) (4.715,-2.280) (4.750,-2.285) (4.786,-2.289) (4.822,-2.292) (4.858,-2.295) (4.894,-2.297) (4.930,-2.299) (4.950,-2.300)}
\LC{(1.331,0.166) (1.365,0.153) (1.398,0.140) (1.432,0.127) (1.465,0.114) (1.499,0.101) (1.533,0.088) (1.566,0.076) (1.600,0.063) (1.633,0.049) (1.667,0.036) (1.700,0.023) (1.734,0.010) (1.767,-0.003) (1.801,-0.017) (1.834,-0.030) (1.867,-0.044) (1.901,-0.057) (1.934,-0.071) (1.967,-0.085) (2.000,-0.099) (2.033,-0.113) (2.066,-0.128) (2.099,-0.142) (2.132,-0.157) (2.165,-0.171) (2.198,-0.186) (2.231,-0.201) (2.263,-0.217) (2.296,-0.232) (2.328,-0.248) (2.361,-0.263) (2.393,-0.279) (2.425,-0.295) (2.458,-0.311) (2.490,-0.327) (2.522,-0.343) (2.554,-0.360) (2.586,-0.376) (2.618,-0.393) (2.650,-0.409) (2.682,-0.426) (2.714,-0.442) (2.746,-0.459) (2.778,-0.476) (2.809,-0.493) (2.841,-0.510) (2.873,-0.526) (2.905,-0.543) (2.937,-0.560) (2.969,-0.577) (3.000,-0.593) (3.032,-0.610) (3.064,-0.627) (3.096,-0.643) (3.128,-0.660) (3.160,-0.676) (3.193,-0.692) (3.225,-0.708) (3.257,-0.724) (3.289,-0.740) (3.322,-0.756) (3.354,-0.772) (3.387,-0.787) (3.419,-0.803) (3.452,-0.818) (3.484,-0.833) (3.517,-0.848) (3.550,-0.862) (3.583,-0.876) (3.616,-0.891) (3.650,-0.905) (3.683,-0.918) (3.716,-0.932) (3.750,-0.945) (3.783,-0.958) (3.817,-0.971) (3.851,-0.983) (3.885,-0.995) (3.919,-1.007) (3.953,-1.019) (3.987,-1.030) (4.021,-1.041) (4.055,-1.052) (4.090,-1.062) (4.124,-1.072) (4.159,-1.082) (4.194,-1.092) (4.229,-1.101) (4.264,-1.109) (4.299,-1.118) (4.334,-1.125) (4.369,-1.133) (4.404,-1.140) (4.440,-1.147) (4.475,-1.153) (4.511,-1.159) (4.546,-1.165) (4.582,-1.170) (4.617,-1.175) (4.653,-1.179) (4.689,-1.183) (4.725,-1.187) (4.761,-1.190) (4.797,-1.193) (4.832,-1.195) (4.868,-1.197) (4.904,-1.198) (4.940,-1.199) (4.950,-1.199)}
\LC{(1.362,0.297) (1.398,0.294) (1.434,0.291) (1.470,0.288) (1.506,0.286) (1.542,0.283) (1.578,0.281) (1.614,0.279) (1.650,0.278) (1.686,0.276) (1.722,0.275) (1.757,0.273) (1.793,0.272) (1.829,0.271) (1.865,0.270) (1.901,0.269) (1.937,0.268) (1.973,0.267) (2.009,0.266) (2.045,0.265) (2.081,0.264) (2.117,0.262) (2.153,0.261) (2.189,0.259) (2.225,0.257) (2.261,0.255) (2.297,0.252) (2.333,0.249) (2.369,0.246) (2.405,0.243) (2.440,0.239) (2.476,0.234) (2.512,0.229) (2.547,0.224) (2.583,0.218) (2.618,0.212) (2.654,0.205) (2.689,0.198) (2.724,0.190) (2.759,0.182) (2.794,0.173) (2.829,0.164) (2.864,0.155) (2.898,0.145) (2.933,0.135) (2.968,0.125) (3.002,0.114) (3.036,0.104) (3.071,0.093) (3.105,0.081) (3.139,0.070) (3.173,0.059) (3.207,0.047) (3.241,0.035) (3.275,0.024) (3.309,0.012) (3.343,0.000) (3.377,-0.011) (3.412,-0.023) (3.446,-0.035) (3.480,-0.046) (3.514,-0.057) (3.548,-0.069) (3.582,-0.080) (3.617,-0.091) (3.651,-0.102) (3.685,-0.113) (3.720,-0.123) (3.754,-0.134) (3.789,-0.144) (3.823,-0.154) (3.858,-0.164) (3.892,-0.173) (3.927,-0.183) (3.962,-0.192) (3.997,-0.201) (4.032,-0.209) (4.067,-0.218) (4.102,-0.226) (4.137,-0.233) (4.172,-0.241) (4.208,-0.248) (4.243,-0.255) (4.278,-0.262) (4.314,-0.268) (4.349,-0.274) (4.385,-0.280) (4.420,-0.286) (4.456,-0.291) (4.492,-0.296) (4.527,-0.300) (4.563,-0.305) (4.599,-0.308) (4.635,-0.312) (4.671,-0.315) (4.706,-0.318) (4.742,-0.321) (4.778,-0.323) (4.814,-0.325) (4.850,-0.326) (4.886,-0.328) (4.922,-0.329) (4.950,-0.329)}
\LC{(3.366,1.285) (3.373,1.250) (3.384,1.215) (3.399,1.183) (3.417,1.152) (3.438,1.123) (3.462,1.095) (3.487,1.070) (3.514,1.046) (3.543,1.024) (3.572,1.004) (3.603,0.985) (3.634,0.967) (3.666,0.951) (3.699,0.935) (3.732,0.921) (3.765,0.907) (3.799,0.895) (3.833,0.883) (3.867,0.872) (3.902,0.861) (3.936,0.851) (3.971,0.842) (4.006,0.834) (4.041,0.825) (4.076,0.818) (4.111,0.811) (4.147,0.804) (4.182,0.797) (4.218,0.791) (4.253,0.786) (4.289,0.781) (4.325,0.776) (4.360,0.771) (4.396,0.767) (4.432,0.763) (4.468,0.759) (4.503,0.756) (4.539,0.753) (4.575,0.750) (4.611,0.747) (4.647,0.745) (4.683,0.743) (4.719,0.741) (4.755,0.739) (4.791,0.738) (4.827,0.737) (4.863,0.736) (4.899,0.735) (4.935,0.735) (4.950,0.734)}
\LC{(3.585,1.402) (3.621,1.398) (3.657,1.394) (3.692,1.389) (3.728,1.385) (3.764,1.381) (3.800,1.377) (3.835,1.374) (3.871,1.371) (3.907,1.369) (3.943,1.367) (3.979,1.365) (4.015,1.364) (4.051,1.363) (4.087,1.362) (4.123,1.362) (4.159,1.362) (4.195,1.362) (4.231,1.362) (4.267,1.362) (4.303,1.362) (4.339,1.363) (4.375,1.363) (4.411,1.364) (4.447,1.364) (4.483,1.365) (4.519,1.366) (4.555,1.366) (4.591,1.367) (4.627,1.367) (4.663,1.368) (4.699,1.368) (4.735,1.369) (4.771,1.369) (4.807,1.370) (4.843,1.370) (4.879,1.370) (4.915,1.370) (4.950,1.371)}
\LC{(3.585,1.615) (3.620,1.621) (3.655,1.630) (3.690,1.640) (3.724,1.650) (3.759,1.661) (3.793,1.672) (3.827,1.683) (3.862,1.694) (3.896,1.704) (3.931,1.715) (3.965,1.725) (4.000,1.735) (4.034,1.745) (4.069,1.755) (4.104,1.764) (4.138,1.773) (4.173,1.782) (4.208,1.790) (4.243,1.799) (4.279,1.807) (4.314,1.814) (4.349,1.821) (4.384,1.828) (4.420,1.835) (4.455,1.841) (4.491,1.847) (4.526,1.853) (4.562,1.858) (4.598,1.862) (4.633,1.867) (4.669,1.871) (4.705,1.874) (4.741,1.878) (4.777,1.880) (4.813,1.883) (4.848,1.885) (4.884,1.886) (4.920,1.887) (4.950,1.888)}
\LC{(3.422,1.715) (3.428,1.751) (3.437,1.785) (3.449,1.819) (3.463,1.852) (3.479,1.884) (3.497,1.916) (3.516,1.946) (3.537,1.976) (3.558,2.005) (3.581,2.033) (3.604,2.060) (3.628,2.087) (3.652,2.114) (3.677,2.140) (3.703,2.165) (3.729,2.190) (3.755,2.214) (3.782,2.238) (3.809,2.262) (3.837,2.285) (3.865,2.307) (3.893,2.329) (3.922,2.351) (3.951,2.372) (3.981,2.393) (4.011,2.413) (4.041,2.433) (4.071,2.452) (4.102,2.471) (4.133,2.489) (4.164,2.507) (4.196,2.524) (4.228,2.541) (4.260,2.557) (4.292,2.573) (4.325,2.588) (4.358,2.602) (4.391,2.616) (4.425,2.629) (4.458,2.642) (4.492,2.654) (4.527,2.665) (4.561,2.675) (4.596,2.685) (4.630,2.694) (4.665,2.703) (4.701,2.710) (4.736,2.717) (4.772,2.723) (4.807,2.728) (4.843,2.732) (4.879,2.735) (4.915,2.738) (4.950,2.739)}
\LC{(3.069,1.715) (3.071,1.751) (3.073,1.787) (3.077,1.823) (3.081,1.858) (3.087,1.894) (3.093,1.929) (3.100,1.965) (3.108,2.000) (3.117,2.035) (3.126,2.069) (3.137,2.104) (3.148,2.138) (3.160,2.172) (3.172,2.206) (3.185,2.240) (3.198,2.273) (3.212,2.306) (3.227,2.339) (3.242,2.372) (3.257,2.404) (3.273,2.437) (3.289,2.469) (3.306,2.501) (3.323,2.532) (3.341,2.564) (3.358,2.595) (3.376,2.626) (3.395,2.657) (3.413,2.688) (3.432,2.719) (3.452,2.749) (3.471,2.779) (3.491,2.810) (3.511,2.840) (3.531,2.869) (3.551,2.899) (3.572,2.929) (3.592,2.958) (3.613,2.987) (3.634,3.017) (3.656,3.046) (3.677,3.075) (3.699,3.103) (3.720,3.132) (3.742,3.161) (3.764,3.189) (3.786,3.218) (3.809,3.246) (3.831,3.274) (3.854,3.302) (3.876,3.330) (3.899,3.358) (3.922,3.386) (3.945,3.413) (3.968,3.441) (3.991,3.469) (4.014,3.496) (4.038,3.524) (4.061,3.551) (4.085,3.578) (4.108,3.605) (4.132,3.633) (4.155,3.660) (4.179,3.687) (4.203,3.714) (4.227,3.741) (4.250,3.768) (4.274,3.795) (4.298,3.822) (4.322,3.849) (4.346,3.876) (4.370,3.903) (4.394,3.930) (4.417,3.957) (4.441,3.984) (4.465,4.011) (4.488,4.038) (4.512,4.065) (4.535,4.093) (4.558,4.120) (4.581,4.148) (4.604,4.176) (4.626,4.204) (4.648,4.233) (4.669,4.262) (4.689,4.292) (4.709,4.322) (4.727,4.354) (4.742,4.386) (4.755,4.420) (4.764,4.450)}
\LC{(2.580,1.715) (2.581,1.751) (2.582,1.787) (2.583,1.823) (2.585,1.859) (2.587,1.895) (2.590,1.931) (2.593,1.967) (2.596,2.002) (2.600,2.038) (2.604,2.074) (2.608,2.110) (2.613,2.146) (2.618,2.181) (2.623,2.217) (2.629,2.252) (2.634,2.288) (2.641,2.323) (2.647,2.359) (2.654,2.394) (2.661,2.429) (2.668,2.465) (2.676,2.500) (2.683,2.535) (2.691,2.570) (2.700,2.605) (2.708,2.640) (2.716,2.675) (2.725,2.710) (2.734,2.745) (2.743,2.780) (2.752,2.815) (2.761,2.850) (2.770,2.884) (2.780,2.919) (2.789,2.954) (2.799,2.989) (2.808,3.023) (2.818,3.058) (2.828,3.093) (2.837,3.127) (2.847,3.162) (2.857,3.197) (2.867,3.231) (2.876,3.266) (2.886,3.301) (2.896,3.335) (2.905,3.370) (2.914,3.405) (2.924,3.440) (2.933,3.474) (2.942,3.509) (2.951,3.544) (2.960,3.579) (2.969,3.614) (2.977,3.649) (2.985,3.684) (2.993,3.719) (3.001,3.754) (3.009,3.789) (3.016,3.825) (3.023,3.860) (3.030,3.895) (3.037,3.931) (3.043,3.966) (3.049,4.002) (3.054,4.037) (3.059,4.073) (3.064,4.108) (3.069,4.144) (3.073,4.180) (3.076,4.216) (3.080,4.252) (3.082,4.287) (3.085,4.323) (3.087,4.359) (3.088,4.395) (3.089,4.431) (3.090,4.450)}
\LC{(1.996,1.715) (1.996,1.751) (1.996,1.787) (1.997,1.823) (1.997,1.859) (1.998,1.895) (1.999,1.931) (2.000,1.967) (2.001,2.003) (2.002,2.039) (2.003,2.075) (2.005,2.111) (2.007,2.147) (2.008,2.183) (2.010,2.219) (2.012,2.255) (2.014,2.291) (2.016,2.327) (2.019,2.362) (2.021,2.398) (2.023,2.434) (2.026,2.470) (2.029,2.506) (2.031,2.542) (2.034,2.578) (2.037,2.614) (2.040,2.650) (2.043,2.686) (2.046,2.721) (2.049,2.757) (2.052,2.793) (2.056,2.829) (2.059,2.865) (2.062,2.901) (2.066,2.937) (2.069,2.972) (2.072,3.008) (2.076,3.044) (2.079,3.080) (2.083,3.116) (2.086,3.152) (2.090,3.187) (2.093,3.223) (2.097,3.259) (2.100,3.295) (2.104,3.331) (2.107,3.367) (2.110,3.402) (2.114,3.438) (2.117,3.474) (2.120,3.510) (2.123,3.546) (2.126,3.582) (2.129,3.618) (2.132,3.653) (2.135,3.689) (2.138,3.725) (2.141,3.761) (2.143,3.797) (2.146,3.833) (2.148,3.869) (2.151,3.905) (2.153,3.941) (2.155,3.977) (2.157,4.013) (2.159,4.049) (2.161,4.084) (2.162,4.120) (2.164,4.156) (2.165,4.192) (2.166,4.228) (2.167,4.264) (2.168,4.300) (2.169,4.336) (2.169,4.372) (2.170,4.408) (2.170,4.444) (2.170,4.450)}
\LC{(1.362,1.715) (1.363,1.751) (1.363,1.787) (1.363,1.823) (1.363,1.859) (1.363,1.895) (1.363,1.931) (1.363,1.967) (1.363,2.003) (1.362,2.039) (1.362,2.075) (1.362,2.111) (1.362,2.147) (1.361,2.183) (1.361,2.219) (1.360,2.255) (1.360,2.291) (1.359,2.327) (1.358,2.363) (1.358,2.399) (1.357,2.435) (1.356,2.471) (1.355,2.507) (1.354,2.543) (1.354,2.579) (1.353,2.615) (1.352,2.651) (1.351,2.687) (1.350,2.723) (1.349,2.759) (1.349,2.795) (1.348,2.831) (1.347,2.867) (1.346,2.903) (1.346,2.939) (1.345,2.975) (1.344,3.011) (1.344,3.047) (1.343,3.083) (1.342,3.119) (1.342,3.155) (1.341,3.191) (1.341,3.227) (1.340,3.263) (1.340,3.299) (1.340,3.335) (1.339,3.371) (1.339,3.407) (1.339,3.443) (1.338,3.479) (1.338,3.515) (1.338,3.551) (1.338,3.587) (1.338,3.623) (1.338,3.659) (1.337,3.695) (1.337,3.731) (1.337,3.767) (1.337,3.803) (1.337,3.839) (1.337,3.875) (1.337,3.911) (1.337,3.947) (1.337,3.983) (1.337,4.019) (1.337,4.055) (1.337,4.091) (1.337,4.127) (1.337,4.163) (1.337,4.199) (1.337,4.235) (1.337,4.271) (1.337,4.307) (1.337,4.343) (1.337,4.379) (1.337,4.415) (1.337,4.450)}

    % Disco cargado, resto neutros - flechas
% Generado por scripts/generar_datos.py (no editar a mano).
\FL{-4.553}{-3.230}{-165.4}
\FL{-4.383}{-3.753}{-140.8}
\FL{-3.830}{-3.980}{-115.9}
\FL{-3.424}{-4.120}{-105.9}
\FL{-3.106}{-4.217}{-91.9}
\FL{-2.703}{-4.114}{-73.5}
\FL{-2.302}{-3.932}{-68.0}
\FL{-1.713}{-3.653}{-68.7}
\FL{-0.860}{-3.200}{-77.3}
\FL{-3.857}{-1.417}{156.6}
\FL{-4.287}{-2.170}{156.6}
\FL{-4.533}{-2.677}{168.3}
\FL{-0.860}{3.677}{91.2}
\FL{-1.746}{3.676}{93.1}
\FL{-2.673}{3.665}{105.1}
\FL{-3.862}{3.071}{149.4}
\FL{-3.876}{1.729}{175.8}
\FL{3.149}{0.750}{-24.5}
\FL{1.791}{0.834}{45.2}
\FL{1.562}{0.945}{58.1}
\FL{1.392}{1.001}{65.0}
\FL{1.264}{1.050}{70.2}
\FL{1.136}{1.074}{72.6}
\FL{1.017}{1.098}{74.1}
\FL{0.896}{1.150}{76.3}
\FL{0.923}{2.635}{96.1}
\FL{0.027}{2.594}{114.5}
\FL{0.097}{1.561}{129.1}
\FL{-0.022}{1.224}{140.2}
\FL{-0.109}{0.963}{150.9}
\FL{-0.375}{0.761}{162.7}
\FL{-2.260}{0.941}{-176.9}
\FL{-0.366}{0.069}{-134.0}
\FL{-0.070}{0.054}{-127.6}
\FL{0.093}{0.022}{-126.0}
\FL{0.206}{-0.010}{-129.4}
\FL{0.387}{-0.209}{-150.3}
\FL{0.394}{-0.321}{-156.3}
\FL{0.410}{-0.430}{-155.5}
\FL{0.417}{-0.649}{-152.0}
\FL{0.323}{-1.091}{-146.4}
\FL{0.554}{-2.271}{-99.5}
\FL{1.531}{-2.275}{-80.0}
\FL{2.470}{-2.191}{-59.5}
\FL{2.976}{-1.384}{-39.3}
\FL{3.107}{-0.649}{-27.2}
\FL{3.173}{0.059}{-18.6}
\FL{4.076}{0.818}{-11.9}
\FL{4.267}{1.362}{0.5}
\FL{4.267}{1.804}{12.6}
\FL{4.041}{2.433}{32.8}
\FL{3.757}{3.180}{52.2}
\FL{2.828}{3.093}{74.4}
\FL{2.081}{3.092}{84.5}
\FL{1.343}{3.083}{90.9}

  \end{scope}
  \foreach \f in {\formaA,\formaB,\formaD}{\fill[top color=azulvelo!80!white, bottom color=metal, draw=azulprof, line width=0.55pt] \f;}
  \fill[top color=azulcielo, bottom color=azulmedio, draw=azulprof, line width=0.55pt] \formaC;
  % Disco cargado, resto neutros - cargas
% Generado por scripts/generar_datos.py (no editar a mano).
\CP{-3.977}{-3.386}
\CP{-3.165}{-4.063}
\CP{-2.068}{-3.062}
\CP{-3.817}{-2.494}
\CN{0.084}{-0.581}
\CN{0.336}{-0.291}
\CN{-0.150}{-0.221}
\CP{-1.508}{3.010}
\CP{-2.910}{2.402}
\CN{-0.810}{1.090}
\CN{-0.490}{1.714}
\CP{1.447}{0.581}
\CP{1.197}{0.888}
\CP{0.865}{0.990}
\CP{0.514}{0.895}
\CP{0.258}{0.593}
\CP{0.216}{0.264}
\CP{0.314}{-0.000}
\CP{0.476}{-0.169}
\CP{0.752}{-0.282}
\CP{1.142}{-0.219}
\CP{1.448}{0.123}
\CP{3.710}{1.369}
\CP{3.270}{1.840}
\CP{1.891}{1.840}
\CN{1.170}{1.160}
\CN{1.531}{1.160}
\CN{2.452}{1.160}

  \draw[azulprof, line width=0.7pt] (-5,-4.5) rectangle (5,4.5);
  % rótulos
  \node[font=\footnotesize, fill=white, inner sep=1.5pt, rounded corners=1pt, align=center] at (-3.55,-1.55) {A\\[-1pt]\scriptsize$\phi=\csname philibreA\endcsname$};
  \node[font=\footnotesize, inner sep=1pt, align=center] at (-1.7,2.05) {B\\[-1pt]\scriptsize$\phi=\csname philibreB\endcsname$};
  \node[font=\footnotesize, fill=white, inner sep=1.5pt, rounded corners=1pt, align=center] at (2.05,-0.6) {C\ ($Q>0$)\\[-1pt]\scriptsize$\phi=\csname philibreC\endcsname$};
  \node[font=\footnotesize, fill=white, inner sep=1.5pt, rounded corners=1pt, align=center] at (2.3,2.45) {D\\[-1pt]\scriptsize$\phi=\csname philibreD\endcsname$};
  \node[font=\scriptsize, text=grispizarra, anchor=north east] at (4.9,-3.8) {caja a tierra, $\phi=0$};
\end{tikzpicture}
\caption{Estado «un conductor cargado, los demás neutros», que define la columna $\mathrm{C}$ de los coeficientes de potencial (sección de conductores
cilíndricos, calculado numéricamente; potenciales normalizados a $\phi_{\mathrm{C}}=1$). Las líneas salen del disco cargado; los conductores neutros
reciben líneas por la cara próxima, donde se induce carga negativa, y las reemiten por la opuesta. Todos los potenciales son positivos y el del
conductor cargado es el mayor: $\phi_{\mathrm{C}}>\phi_{i}>0$.}
\label{fig:sistema-libre}
\end{figure}
